\documentclass[11pt]{amsart}
\usepackage{amssymb,amsmath,amsthm,mathtools}
\usepackage[margin=1.1in]{geometry}
\usepackage{hyperref}

\newtheorem{theorem}{Theorem}
\newtheorem{corollary}[theorem]{Corollary}
\newtheorem{proposition}[theorem]{Proposition}
\newtheorem{lemma}[theorem]{Lemma}
\theoremstyle{definition}
\newtheorem{definition}[theorem]{Definition}
\newtheorem{remark}[theorem]{Remark}
\newtheorem{statement}[theorem]{Statement}

\newcommand{\Qb}{\overline{\mathbb{Q}}}
\newcommand{\C}{\mathbb{C}}
\newcommand{\Q}{\mathbb{Q}}
\newcommand{\Z}{\mathbb{Z}}
\newcommand{\R}{\mathbb{R}}
\newcommand{\Ga}{\mathbb{G}_{\mathrm a}}
\newcommand{\Gm}{\mathbb{G}_{\mathrm m}}
\newcommand{\Lg}{\mathcal{L}}
\newcommand{\Lt}{\widetilde{\mathcal{L}}}
\newcommand{\vp}{\varpi}

\title[Four exponentials at CM periods]{Four exponentials at CM periods, and a bilinear statement\\ for two CM elliptic curves}
\author{Draft --- not for circulation}
\date{30 October 2026 (draft)}

\begin{document}

\begin{abstract}
Let $\Omega$ be a nonzero period of an elliptic curve with complex multiplication, defined over $\Qb$. We prove the four
exponentials conjecture at $t=c\,\Omega$ and at $t=c\,\Omega/\pi$ (and hence at their inverses) for every $c\in\Qb^\times$:
for instance, $2^{\vp}$ and
$3^{\vp}$ are not both algebraic, and neither are $2^{\vp/\pi}$ and $3^{\vp/\pi}$, where $\vp$ is the real period of
$y^2=4x^3-4x$. We also show that one of $e^{\pi^2}$, $e^{\pi\Omega}$, $e^{\pi^2/\Omega}$ is transcendental. These follow from
dimension bounds of four-exponentials shape on $\Gm^2\times E$ and $\Gm^2\times E^\natural$. A second application of the
same mechanism, to two non-isogenous CM curves, gives a bilinear statement: if $u$ is a non-torsion elliptic logarithm
on $y^2=4x^3-4x$, $\Omega_\rho$ a period of $y^2=4x^3-4$, and $a,t$ are algebraic, then
$\vp\,\Omega_\rho\,(\zeta(u)+t)/(\pi(u+a))$ is transcendental. The proofs of these results apply Waldschmidt's algebraic
subgroup theorem with period term (1988), in the $K$-version of Roy and Waldschmidt, directly to explicit quotient groups
whose connected algebraic subgroups we list. The period term comes from a period of $\Gm\times E$ (resp.\ $\Gm\times
E^\natural$) lying on the relevant line, and in the two-curve case from the Legendre relations of both curves. We state precisely
what the proofs use from the literature, and write out the one step of Waldschmidt's auxiliary-function construction
(multiplicities along a subspace) that is only indicated in print. We also record a single-curve theorem
(``Theorem M'') with a second proof by Laurent's identity and W\"ustholz's analytic subgroup theorem.
\end{abstract}

\maketitle

\section{Introduction}

Let $\Qb$ be the algebraic closure of $\Q$ in $\C$, $\Lg=\{z\in\C: e^z\in\Qb^\times\}$ and $\Lt=\Qb+\Qb\Lg$. For an
elliptic curve $E\colon y^2=4x^3-g_2x-g_3$ with $g_2,g_3\in\Qb$ we write $\wp,\zeta$ for its Weierstrass functions,
$\Lambda_E$ for its period lattice, $\eta(\lambda)=\zeta(z+\lambda)-\zeta(z)$ for the quasi-periods, and
$\Lg_E=\{u\in\C: u\in\Lambda_E \text{ or } \wp(u)\in\Qb\}$. We use two curves throughout:
\begin{align*}
E_i&\colon y^2=4x^3-4x, & \Lambda&=\vp\,\Z[i], & \vp&=\tfrac{\Gamma(1/4)^2}{2\sqrt{2\pi}}, & \eta(\vp)&=\tfrac{\pi}{\vp};\\
E_\rho&\colon y^2=4x^3-4, & \Lambda&=\Omega_\rho\Z[\rho], & \Omega_\rho&=\tfrac{\Gamma(1/3)^3}{2^{4/3}\pi}, & \eta(\Omega_\rho)&=\tfrac{2\pi}{\sqrt3\,\Omega_\rho}.
\end{align*}
(The quasi-periods follow from the Legendre relation and $\zeta(iz)=-i\zeta(z)$, $\zeta(\rho z)=\bar\rho\zeta(z)$. The real
periods are Beta integrals, $\vp=\int_1^\infty(x^3-x)^{-1/2}dx=\tfrac12B(\tfrac14,\tfrac12)$ and
$\Omega_\rho=\int_1^\infty(x^3-1)^{-1/2}dx=\tfrac13B(\tfrac16,\tfrac12)$, and the $\Gamma$-forms follow from the reflection and
duplication formulas; these are instances of the Chowla--Selberg formula \cite{Gr78}.)

For a CM curve $E$ let $E^\natural$ be its universal vector extension and $\Lambda(E^\natural)$ the logarithms of its
algebraic points, written in coordinates $(z,w)$ (base, and conjugate-eigen fibre; \S\ref{sec:MX}). For $E_i$, $E_\rho$
the fibre coordinate can be normalised so that $X=(u,\zeta(u))\in\Lambda(E^\natural)$ for every $u\in\Lg_E\setminus\Lambda_E$,
and in general $w(X)-\zeta^*(z(X))\in\Qb$ for $X\in\Lambda(E^\natural)$ with $z(X)\notin\Lambda_E$, where
$\zeta^*=\zeta-\kappa_1\,\mathrm{id}$ and $\eta^*=\eta-\kappa_1\,\mathrm{id}$ are the conjugate-eigen versions of $\zeta$ and $\eta$
($\kappa_1\in\Qb$, $\kappa_1=0$ for $E_i,E_\rho$; \S\ref{sec:MX}).

\subsection*{Four exponentials at CM periods}
The four exponentials conjecture asserts that if $\log\alpha,\log\beta$ are linearly independent over $\Q$ and
$t\notin\Q$, then $\alpha^t=e^{t\log\alpha}$ and $\beta^t=e^{t\log\beta}$ are not both algebraic. It is known for algebraic $t$ (Gel'fond--Schneider)
and in several conjugation-generic situations (Diaz \cite{Di04}; Waldschmidt \cite{W07}). We know of no fixed
transcendental $t$ for which it was known for all pairs, including real pairs such as $(2,3)$ \cite{W22}.

\begin{theorem}\label{thm:A}
Let $E/\Qb$ have complex multiplication, let $\Omega\neq0$ be a period of $E$, $c\in\Qb^\times$ and
$\alpha,\beta\in\Qb^\times$. Fix logarithms $\log\alpha$, $\log\beta$ and put $\alpha^z=e^{z\log\alpha}$,
$\beta^z=e^{z\log\beta}$.
\begin{itemize}
\item[(a)] If $\log\alpha,\log\beta$ are linearly independent over $\Q$, then $\alpha^{c\Omega}$ and $\beta^{c\Omega}$ are
not both algebraic.
\item[(b)] If $\log\alpha,\log\beta$ are linearly independent over $\Q$, then $\alpha^{c\Omega/\pi}$ and
$\beta^{c\Omega/\pi}$ are not both algebraic.
\end{itemize}
\end{theorem}

For example (with the real logarithms), $2^{\vp}$ and $3^{\vp}$ are not both algebraic; nor are $2^{\vp/\pi}$ and $3^{\vp/\pi}$, nor $2^{\Omega_\rho}$
and $3^{\Omega_\rho}$. Since the four exponentials statement is symmetric under $t\leftrightarrow1/t$ (if $\alpha^{1/t}$,
$\beta^{1/t}$ are algebraic, apply the statement at $t$ to them), the same holds at $t\in\Qb^\times/\Omega$ and
$t\in\Qb^\times\pi/\Omega$. Theorem~\ref{thm:A} consists of Corollaries~\ref{cor:MGn2}
and~\ref{cor:MG2} of Theorems~\ref{thm:MGn} and~\ref{thm:MG} (\S\ref{sec:MG}). These are dimension bounds that also
allow elliptic logarithms or fibre values in the second row. Another consequence is:

\begin{corollary}\label{cor:MG1intro}
For every CM curve with period $\Omega$, at least one of $e^{\pi^2}$, $e^{\pi\Omega}$, $e^{\pi^2/\Omega}$ is transcendental.
On $y^2=4x^3-4x$: one of $e^{\pi^2}\approx19333.69$, $e^{\pi\vp}\approx3779.84$, $e^{\pi^2/\vp}\approx43.12$ is
transcendental.
\end{corollary}

Part of Theorem~\ref{thm:A} also follows from the six exponentials theorem and complex conjugation (Diaz \cite{Di04};
Waldschmidt \cite[Cor.~2.7]{W07}): for $t\in\R\cup i\R$ and pairs with $\log\alpha,\log\beta,\overline{\log\beta}$ linearly
independent over $\Q$ (e.g.\ $(2,1+2i)$), and, exchanging the roles, for pairs with $\log\alpha,\log\beta\in\R\cup i\R$ at
any $t$ with $1,t,\bar t$ linearly independent over $\Q$. Since every $t$ in Theorem~\ref{thm:A} is an algebraic
multiple of a real number, these arguments together already cover every pair at $t\notin\R\cup i\R$ (see the overlap
paragraph after Corollary~\ref{cor:MG2}). They do not apply to real or purely imaginary $t$ --- such as $t=\vp$,
$\vp/\pi$, $\Omega_\rho$ --- with pairs whose logarithms span a conjugation-stable space, such as $(2,3)$, which is where
our examples lie, nor to Corollary~\ref{cor:MG1intro}.

\subsection*{Two CM curves}
The same mechanism, applied to two non-isogenous CM curves, gives a statement with no logarithm of $\Gm$ at all (and
a twisted variant, Theorem~\ref{thm:MXP} in \S\ref{sec:MXP}). Applied to periods, it also gives two statements of degree
two about periods of the two curves (Corollary~\ref{cor:deg2}): $\pi,\vp,\Omega_\rho,\vp\Omega_\rho$ are linearly independent
over $\Qb$, and so are $1,\vp,\Omega_\rho,\vp\Omega_\rho/\pi$. These are not consequences of the analytic subgroup theorem,
and we flag them as the most delicate claims of the paper.

\begin{theorem}[Theorem MX]\label{thm:MX}
Let $X\in\Lambda(E_i^\natural)$, $Y\in\Lambda(E_\rho^\natural)$, $a,t\in\Qb$ and $\beta\in\Qb^\times$. If
\[
\bigl(z(X)+z(Y)+a\bigr)\cdot\beta\,\eta(\Omega_\rho)\;=\;\bigl(w(X)+w(Y)+t\bigr)\cdot\vp ,
\]
then $a=0$, $z(X)\in\Q(i)\vp$ and $z(Y)\in\Q(\rho)\Omega_\rho$; in particular $\exp_{E_i}(z(X))$ and $\exp_{E_\rho}(z(Y))$ are
torsion points. In particular, for $u\in\Lg_{E_i}\setminus\Lambda$ and $u'\in\Lg_{E_\rho}\setminus\Lambda_\rho$ not both torsion,
$(u+u'+a)\,\beta\,\eta(\Omega_\rho)\neq(\zeta_i(u)+\zeta_\rho(u')+t)\,\vp$.
An analogous statement holds for any two non-isogenous CM elliptic curves (Remark~\ref{rem:general}).
\end{theorem}

\begin{corollary}\label{cor:MX1}
For $u\in\Lg_{E_i}\setminus\Q(i)\vp$ and $a,t\in\Qb$,
\[
\frac{\vp\,\Omega_\rho\,(\zeta_i(u)+t)}{\pi\,(u+a)}\notin\Qb,\qquad\text{equivalently}\qquad
\frac{\zeta_i(u)+t}{u+a}\notin \Qb\cdot\frac{\pi^{5/2}}{\Gamma(1/4)^2\,\Gamma(1/3)^3}.
\]
Symmetrically for non-torsion $u'\in\Lg_{E_\rho}$ with $\zeta_\rho$.
\end{corollary}

For example, for the point with $\wp(u)=2$ one has $u=0.72694593546890819853\ldots$ and
$\vp\Omega_\rho\zeta(u)/(\pi u)=3.76379062156897182773\ldots$ (computed independently twice, by theta functions and by
the Laurent series of $\zeta$; an illustration only). The case $X=Y=0$ of Theorem~\ref{thm:MX} gives:

\begin{corollary}\label{cor:gamma}
$\Gamma(1/4)^2\,\Gamma(1/3)^3/\pi^{5/2}$ is transcendental.
\end{corollary}

This is immediate from W\"ustholz's theorem (Remark~\ref{rem:gamma}), and probably also from the Wolfart--W\"ustholz theorem
on Beta values \cite{WW85} (not checked); we found it stated nowhere, but do not claim it as new. Algebraic independence of $\pi,\Gamma(1/4),\Gamma(1/3)$ is open, so Corollary~\ref{cor:MX1} does not follow from
known algebraic independence results.

\medskip
The single-curve version of the same count gives the following theorem, which turns out to be classical in substance.

\begin{theorem}[Theorem M]\label{thm:M}
Let $E/\Qb$ have complex multiplication by $k$, and let $\Omega\neq0$ be a period. If $u\in\Qb\Lg_E$ and $\ell\in\Lt$
satisfy $i\pi\, u=\Omega\,\ell$, then $u=s\Omega$ and $\ell=s\,i\pi$ for some $s\in\Qb$; if $u\in\Lg_E$ then $s\in k$.
\end{theorem}

\begin{corollary}\label{cor:M}
On $E_i$ (and analogously on every CM curve):
(M1) $\wp(c\vp/\pi)\notin\Qb$ for $c\in\Qb^\times$;
(M2) $\wp(\vp\log\alpha/\pi)\notin\Qb$ for $\alpha\in\Qb^\times$ not a root of unity;
(M3) at least one of $e^{\pi^2}$, $\wp(\pi\vp)$ is transcendental; more precisely, if $e^{\pi^2}\in\Qb$ then
$\wp(c\pi\Omega)\notin\Qb$ for every $c\in\Qb^\times$ and every period $\Omega$ of every CM curve;
(M3$'$) for $\alpha\in\Qb\setminus\{0,1\}$ and any branch, one of $e^{i\pi\log\alpha}$, $\wp(\vp\log\alpha)$ is transcendental;
(M$\eta$) $\eta(\vp)\,u\notin\Lt$ for $u\in\Lg_{E_i}\setminus\Q(i)\vp$.
\end{corollary}

\noindent\emph{Novelty, honestly.} With $\ell\in\Lg$, Theorem~\ref{thm:M} is the CM case, with one row of Waldschmidt's matrix a period of
$E\times\Gm$, of his Conjecture~42 in \cite{W06} (four exponentials for $E\times\Gm$). Its special case
$\ell\in\Qb+\Qb\cdot\pi i$, $u\in\Lg_E$, follows in one line from Laurent's \cite[Cor.~3, p.~146]{La82}: for any CM
lattice, a period $\omega\ne0$ and $u\notin\Q\Lambda_E$ with $\wp(u)\in\Qb$, the numbers
$1,\omega,\eta(\omega),2\pi i,\omega H(u),2\pi iu/\omega$ are $\Qb$-linearly independent, where $H=\zeta-\gamma z$ and
$\gamma=0$ when $g_2g_3=0$. (M1) is one line from it ($2\pi iu/\Omega=2ic$), and so is (M$\eta$) modulo
$\Qb+\Qb\vp+\Qb\eta(\vp)+\Qb\pi i$. The remaining items of Corollary~\ref{cor:M} (which only use $u\in\Lg_E$) follow quickly from
Laurent's identity and the analytic subgroup theorem (\S\ref{sec:M2}); we found none of them stated in print (Abe \cite{Abe20} claimed (M1); the preprint was
withdrawn). Masser's theorem \cite{Ma76} ($1,\omega,\eta,2\pi i,\log\sigma$ are $\Qb$-linearly independent on a CM curve, for
$\sigma\in\Qb^\times$ not a root of unity) is
linear and gives none of these; nor do the algebraic independence results of Brownawell--Kubota \cite{BK77}, which have
six-exponentials strength, or the third-kind results surveyed in \cite{W19} (Wüstholz's linear independence of
$1,\omega,\eta,\lambda(u_i,\omega)$). Theorem~\ref{thm:A}, Corollary~\ref{cor:MG1intro}, Theorem~\ref{thm:MX} and
Corollary~\ref{cor:MX1} we found nowhere; the four exponentials problem is listed as open, with no known case of this
kind, in \cite[p.~27]{SMF84} and \cite{W22}. For a recent unrefereed preprint, partly Lean-checked relative to stated
classical inputs, on CM periods of a different kind (irrationality exponent $2$ for real periods of CM curves over $\Q$, and $\mu(\vp^n)=O(n\log 2n)$), see \cite{ChCh26}. For Theorem~\ref{thm:A} the closest published result is \cite[Thm.~4.2]{W05} (\S\ref{sec:MG}). The nearest published statement we know
is Waldschmidt's \cite[Thm.~4.3]{W05}: a bilinear statement for two elliptic curves (non-isogenous, or equal with
$k$-independent first column), but for $2\times9$ matrices
of elliptic logarithms, without quasi-periods and without a $2\times2$ case. His ``Further variations''
\cite{W05b} concerns $2\times2$ minors with entries in $\Lt$ only (no elliptic functions). Laurent \cite{La82} never
mixes two curves; Corollary~\ref{cor:MX1} is the cross-curve analogue of his Cor.~3, with $\eta(\Omega_\rho)$ in place
of $\eta(\vp)$. Sources we could not consult are listed in \S\ref{sec:status}.

\medskip
\noindent\emph{Why MX is different.} After multiplying by $\pi$, every relation in Theorem~\ref{thm:M} with $u\in\Lg_E$ becomes a
$\Qb$-linear relation among periods of $1$-motives, because of Laurent's identity (\ref{eq:laurent}); this is why the
analytic subgroup theorem proves it too. The hypothesis of Theorem~\ref{thm:MX} instead contains the products
$\Omega_\rho\,\zeta_i(u)$ and $\eta(\vp)\,u'$, each pairing a period of one curve with a logarithm of the other. No CM
identity linearises these; by the period conjecture they are genuine degree-two periods (\S\ref{sec:3E}). The extra
strength comes from a second true degree-two period relation (the Legendre relation of the second curve), which the
count sees as a second period on the line.

\section{The subgroup theorem with period term, and the Roy--Waldschmidt count}\label{sec:RW}

Let $G$ be a commutative algebraic group over $\Qb$ of positive dimension, $T_G(\C)$ its tangent space with the
$\Qb$-structure from $G$, $\Omega(G)=\ker\exp_G$ and $\Lambda(G)=\exp_G^{-1}(G(\Qb))$. For a number field $K\subset\C$,
$\Omega_K(G),\Lambda_K(G)$ are their $K$-spans. Let $d_0(G)$ (resp.\ $d_1(G)$) be the largest $d$ with a surjection
$G\to\Ga^d$ (resp.\ $\Gm^d$) over $\Qb$, and $\alpha(G)=2\dim G-2d_0(G)-d_1(G)$. Fix a basis $\beta_1,\dots,\beta_n$ of
$K/\Q$ and $p_\beta\colon T_{G^n}=T_G^n\to T_G$, $(z_j)\mapsto\sum\beta_jz_j$. Let $H$ be the largest connected algebraic
subgroup of $G^n$ defined over $\Qb$ with $p_\beta(T_H)=0$ and put $\alpha_K(G)=\alpha(G^n/H)/n$ \cite[p.~360]{RW93}.
$S_K(G)$ is the set of subspaces $p_\beta(T_{G'})$, $G'\subset G^n$ an algebraic subgroup defined over $\Qb$; they are
$\Qb$-rational.

\begin{theorem}[{\cite[Th\'eor\`eme, p.~363]{RW93}}]\label{thm:RW}
Let $\mathcal E$ be a finite-dimensional $\C$-space with a $\Qb$-structure, $\Psi\colon T_G(\C)\to\mathcal E$ a
$\Qb$-rational linear map, $W\subset V\subset\mathcal E$ with $W$ rational, and $\kappa=\dim_K(V\cap\Psi(\Omega_K(G)))$.
If (i) $V$ contains no nonzero $\Psi(S)$ with $S\in S_K(G)$, then
\[
\dim_K\bigl(V\cap\Psi(\Lambda_K(G))\bigr)\le(\alpha_K(G)-\kappa)\dim_\C(V/W).
\]
\end{theorem}

\noindent By \cite[Remarque 1]{RW93}, (i) holds when $V$ has no nonzero $\Qb$-rational point (and then $W=0$), since every
$\Psi(S)$ is rational. We use only this case. Theorem~\ref{thm:RW} is used below only for the first proof of
Theorem~\ref{thm:M} and in the sketch of \S\ref{sec:MXA}. Theorems~\ref{thm:MG}, \ref{thm:MGn} and~\ref{thm:MX} are proved
from the following theorem directly; the sketch of Theorem~\ref{thm:MXP} and the second derivations of
Remarks~\ref{rem:MG5d}, \ref{rem:MGn5d} and~\ref{rem:count} use its hyperplane consequence, Theorem~\ref{thm:W88k} below.

\begin{theorem}[{\cite[Thm.~4.1, pp.~385--386]{W88}}]\label{thm:W88}
Let $G=\Ga^{d_0}\times\Gm^{d_1}\times G_2$ be a connected commutative algebraic group over $\Qb$ of dimension
$d=d_0+d_1+d_2\ge1$, with projections $\pi_0,\pi_1$ onto the first two factors ($d_0,d_1$ need not be maximal). Let
$V\subset T_G(\C)$ be a subspace of dimension $N<d$, $W\subset V$ a subspace defined over $\Qb$, and $Y\subset V$ a
finitely generated subgroup with $\Gamma=\exp_GY\subset G(\Qb)$; put $\kappa=\mathrm{rk}_\Z(Y\cap\ker\exp_G)$. Then there
is a connected algebraic subgroup $G'\neq G$ defined over $\Qb$ such that, with $\delta=\dim G/G'$,
$\delta_0=\dim\Ga^{d_0}/\pi_0(G')$, $\delta_1=\dim\Gm^{d_1}/\pi_1(G')$, $\delta_2=\delta-\delta_0-\delta_1$,
$\lambda=\mathrm{rk}_\Z\,\Gamma/(\Gamma\cap G')$ and $\tau=\dim W/(W\cap T_{G'})$, one has $\delta>\tau$ and
\[
(\lambda+\delta_1+2\delta_2)(d-N)\le(\delta-\tau)(d_1+2d_2-\kappa).
\]
\end{theorem}

\noindent ($G'=0$ is allowed.) The proof \cite[\S\S6--7]{W88} combines an auxiliary function
\cite[Prop.~6.1]{W88} with Philippon's zero estimate \cite{Ph86}. Prop.~6.1 asks for a polynomial that vanishes with
multiplicity $\ge T$ along $W$ at the points of $\exp_G(Y(S))$, with the saving $\kappa$ in the parameters. The
$\kappa$-saving is proved in \cite[Prop.~2.1 and 2.4]{W83} (an auxiliary function constructed with periods), for $W=0$
and $b=1$ \cite[p.~396]{W88}; for the rest W88 refers to \cite{W83} and to Bertrand \cite[App.~I, Lemme~7 and its
Corollaire, with Lemme~1]{W79} (cited in \cite[p.~396]{W88} as ``Lemma~7''). All our
applications use both $W\ne0$ and $\kappa>0$. We write out the multiplicity step in Appendix~\ref{app:U} (Lemma~\ref{lem:U}).
It holds under the condition $2N\ge d_1+2d_0+\kappa$, which is satisfied in all the applications in \S\S\ref{sec:MG}
and~\ref{sec:MX} (for Theorem~\ref{thm:MG}: $5n-2\ge3n$; Theorem~\ref{thm:MGn}: $5n\ge3n$; Theorem~\ref{thm:MX}:
$4n+2\ge2n+2$; Theorem~\ref{thm:MXP}: $4n+2\ge2n$; Theorem~\ref{thm:MGaff}: $5n+2\ge3n+4$).

\smallskip
We use one piece of linear algebra repeatedly.

\begin{lemma}\label{lem:K}
Let $k\subset K$ be an imaginary quadratic field inside a number field of degree $n$, $\theta_1,\dots,\theta_n$ a basis of
$K/\Q$, and on $k^n$ put $\psi(m)=\sum\theta_jm_j$, $\phi(m)=\sum\theta_j\bar m_j$. Then $(\psi,\phi)\colon k^n\to K\times K$
is a $\Q$-linear bijection; $\ker\psi$ and $\ker\phi$ are $k$-subspaces of dimension $n/2$, and $\ker\psi\cap\ker\phi=0$.
\end{lemma}

\begin{proof}
$k^n\cong K\otimes_\Q k$ via $m\mapsto\sum\theta_j\otimes m_j$, and $K\otimes_\Q k\cong K\times K$ via
$x\otimes y\mapsto(xy,x\bar y)$; the composite is $(\psi,\phi)$. $\psi$ is $k$-linear and $\phi$ is $k$-semilinear, and
both are onto $K$, which has $k$-dimension $n/2$.
\end{proof}

\medskip\noindent\textbf{Where Theorem~\ref{thm:RW} is proved, and the dependence.} The proof \cite[pp.~363--364]{RW93}
has two steps.
\begin{itemize}
\item[(b)] (Reduction, fully written.) With $\Psi_\beta=\Psi\circ p_\beta$, let $H$ be the largest connected subgroup with
$T_H\subset\ker\Psi_\beta$, $G'=G^n/H$, $\Psi'\colon T_{G'}\to\mathcal E$ the induced map, $V'=\Psi'^{-1}(V)$,
$W'=\Psi'^{-1}(W)$. Then $\dim_\Q(V'\cap\Lambda(G'))=n\dim_K(V\cap\Psi(\Lambda_K))$ and
$\dim_\Q(V'\cap\Omega(G'))=n\kappa$, and the theorem follows from case (a) for $G'$ and $\alpha(G')\le n\alpha_K(G)$.
\item[(a)] (Case $K=\Q$, $\Psi=\mathrm{id}$.) \emph{Statement R}: if $V\subset T_G(\C)$ contains no nonzero Lie algebra of
an algebraic subgroup of $G$ defined over $\Qb$, and $W\subset V$ is $\Qb$-rational, then
\[
\dim_\Q(V\cap\Lambda(G))\le\bigl(\alpha(G)-\dim_\Q(V\cap\Omega(G))\bigr)\dim_\C(V/W).
\]
RW93 derive (i)$\Rightarrow$(iii) here ``de mani\`ere purement formelle'' from Theorem~4.1 of Waldschmidt \cite{W88}
(the algebraic subgroup theorem with the period rank $\kappa=\mathrm{rk}(Y\cap\ker\exp_G)$; we checked its statement on
pp.~385--386), by applying Roy's categorical Proposition~2.1 of \cite{R1} to the category $\mathcal G_w$ and the functions
$a',b,c,d,r$ of \cite[\S6]{R1}. Their Remark~2 says that Statement R with the $\kappa$-term was known when
$\dim V=\dim G-1$ (citing \cite[Thm.~1.1]{W88}) or with $\kappa$ replaced by $0$ (\cite[\S6, Remark (i)]{R1}), and is proved
in general ``en raffinant les arguments du \S6 de [R1] (en travaillant avec la fonction $a'$ au lieu de $a''$)''.
\end{itemize}
Step (b) also uses case (a) for $\ker\Psi'$, to see that $\Psi'$ is injective on $\Lambda(G')$; this also follows from
W\"ustholz's analytic subgroup theorem \cite{Wu89}.

\smallskip\noindent\textbf{The dependence, stated precisely.} When $V'=\Psi'^{-1}(V)$ is a \textbf{hyperplane} of
$T_{G'}$ (as in all our applications), Statement R with the $\kappa$-term is in substance the following refinement of
\cite[Thm.~1.1]{W88}, proved in \cite[\S5d, p.~394]{W88}.

\begin{theorem}[{\cite[Thm.~1.1, pp.~375--376, under (5.4) of \S5d, p.~394]{W88}}]\label{thm:W88k}
Let $G=\Ga^{d_0}\times\Gm^{d_1}\times G_2$ be a connected commutative algebraic group over $\Qb$ of dimension $d\ge1$, with
$d_0$, $d_1$ maximal and $d_2=d-d_0-d_1$. Let $V$ be a hyperplane of $T_G(\C)$, $W\subset V$ a subspace of dimension
$t\ge0$ defined over $\Qb$, and $Y=\Z y_1+\dots+\Z y_m\subset V$ a subgroup of rank $m$ with $\exp_GY\subset G(\Qb)$. Put
$\kappa=\mathrm{rk}_\Z(Y\cap\ker\exp_G)$. If
\[
m>(d_1+2d_2-\kappa)(d-1-t),\tag{5.4}
\]
then $V$ contains a nonzero algebraic Lie subalgebra of $T_G(\C)$ defined over $\Qb$, that is, $T_{G'}\subset V$ for some
connected algebraic subgroup $G'\ne0$ of $G$ defined over $\Qb$.
\end{theorem}

\noindent The printed Theorem~1.1 has the hypothesis $m>(d_1+2d_2)(d-1-t)$, without $\kappa$; \cite[\S5d]{W88} proves it
under the weaker hypothesis (5.4) (p.~394 prints $\kappa=\ell-m$, a misprint for $m-\ell$, cf.\ Lemma~5.2 there). The case
$W=0$ ($t=0$) of Theorem~\ref{thm:W88k} is also printed as \cite[Thm.~5.1, p.~352]{W05}, where it is attributed to
\cite[Thm.~1.1]{W88} and to \cite[Thm.~4.1]{W90}. The proof in \cite[\S5d]{W88} applies Theorem~\ref{thm:W88} (in the form
\cite[Cor.~5.3]{W88}) to $G$ with $N=d-1$, and then again to an auxiliary subgroup $G'$ given by \cite[Lemma~5.1]{W88}.
The condition of Lemma~\ref{lem:U} below covers the first application in the cases listed above; we have not checked it
for the second. So every use of Theorem~\ref{thm:W88k} rests on \cite[\S5d]{W88} as published.

For Theorems~\ref{thm:MG}, \ref{thm:MGn} and~\ref{thm:MX} the main proofs bypass Statement R: \S\S\ref{sec:MG}, \ref{sec:MX}
apply Theorem~\ref{thm:W88} to explicit quotient groups and exclude every connected algebraic subgroup. Each of the three
also has a second derivation from Theorem~\ref{thm:W88k}, with a two-line subgroup exclusion (Remarks~\ref{rem:MG5d},
\ref{rem:MGn5d} and~\ref{rem:count}). The sketch of Theorem~\ref{thm:MXP} uses Theorem~\ref{thm:W88k} (a full case check
against Theorem~\ref{thm:W88} is only mentioned there). For Theorem~\ref{thm:M}, the cases (M1), (M3) were checked against
Theorem~\ref{thm:W88} in the same way, and the second proof (\S\ref{sec:M2}) avoids \cite{RW93} for $u\in\Lg_E$; only
$u\in\Qb\Lg_E$ uses Theorem~\ref{thm:RW} as stated.

\medskip\noindent\textbf{Weights used.} $\alpha_K(\Ga)=0$, $\alpha_K(\Gm)=1$, and for an elliptic curve $\alpha_K(E)=1$ if
$E$ has CM by $k$ and $K\supseteq k$, else $2$ \cite[Lemme, p.~361]{RW93}; $\alpha_K$ is additive on products
\cite[(2)]{RW93}. For the universal vector extension $E^\natural$ of a CM curve and $K\supseteq k$ we need:

\begin{lemma}\label{lem:natural}
$\alpha_K(E^\natural)=2+2/n$.
\end{lemma}

\begin{proof}
Use coordinates on $T(E^\natural)=H^1_{\mathrm{dR}}(E)^\vee$ in which $\mathcal O=\mathrm{End}(E)$ acts by its identity
character on the base coordinate and by the conjugate character on the fibre coordinate (the eigenline of
$H^1_{\mathrm{dR}}$ for $\bar a$ is spanned by $x\,dx/y+c\,dx/y$ with $c\in\Qb$, so these coordinates are rational). Let
$H\subset(E^\natural)^n$ be connected with $p_\beta(T_H)=0$, with abelian image $B\subset E^n$. Then $B=B_U$ for a
$k$-subspace $U\subset k^n$ on which both $u\mapsto\sum\beta_ju_j$ (base) and $u\mapsto\sum\beta_j\bar u_j$ (fibre) vanish.
The map $k^n\to K\times K$, $u\mapsto(\sum\beta_ju_j,\sum\beta_j\bar u_j)$ is injective (it is the
composite $k^n\cong K\otimes_\Q k\cong K\times K$, $x\otimes y\mapsto(xy,x\bar y)$), so $B=0$. Hence $H$ lies in the
fibre $\Ga^n$, and the largest such $H$ is the rational hyperplane $\sum\beta_jt_j=0$, of dimension $n-1$. The quotient
has dimension $n+1$ and admits no surjection to $\Ga$ or $\Gm$ (it is a non-split extension of $E^n$ by $\Ga$), so
$\alpha=2(n+1)$ and $\alpha_K=2+2/n$.
\end{proof}

\section{Four exponentials at CM periods}\label{sec:MG}

Throughout this section $E/\Qb$ has complex multiplication by $k$, $\Omega\neq0$ is a period of $E$, and for a number
field $K$ we write $K\Lg$, $K\Lg_E$ for the $K$-spans of $\Lg$ and $\Lg_E$. By Baker's theorem
$K\Lg\cong K\otimes_\Q\Lg$, so elements of $\Lg$ that are $\Q$-linearly independent are $K$-linearly independent.

\begin{theorem}[Theorem MG]\label{thm:MG}
Let $c\in\Qb^\times$ and $s=c\,\Omega/\pi$. For every number field $K\supseteq k$,
\[
\mathcal S_K:=\{x\in K\Lg:\ s\,x\in K\Lg+K\Lg_E\}\quad\text{satisfies}\quad \dim_K\mathcal S_K\le2 ,
\]
and $2\pi i\in\mathcal S_K$ (since $s\cdot2\pi i=2ic\,\Omega\in K\Lg_E$ when $ic\in K$).
\end{theorem}

\begin{proof}
Enlarging $K$ does not change the conclusion (by Baker, $K$-independence in $K\Lg$ is preserved under base change), so
assume $K\supseteq k(i,c)$; $n=[K:\Q]$ is even (the count below closes for every such $n$). Let $G=\Gm\times\Gm\times E$ with coordinates
$(x,y,z)$, $\Psi(x,y,z)=(x,y+z)$, $V=\C(1,s)$. Fix a basis $\theta_1,\dots,\theta_n$ of $K/\Q$, put
$p_\theta\colon T_{G^n}\to T_G$, $(x_j)\mapsto\sum\theta_jx_j$, and $\Phi=\Psi\circ p_\theta$.

\emph{The group.} Connected algebraic subgroups of $\Gm^{2n}\times E^n$ are products $T_L\times B_S$ of a subtorus
(given by a subspace $L\subset\Q^{2n}$) and an abelian subvariety (given by a $k$-subspace $S\subset k^n$). Such a
subgroup has tangent space in $\ker\Phi$ iff $L=0$ and $\psi(S)=0$, where $\psi(m)=\sum\theta_jm_j$; by
Lemma~\ref{lem:K}, $\ker\psi$ has $k$-dimension $n/2$. Put $U=B_{\ker\psi}$ and
$G''=G^n/U=\Gm^{2n}\times A''$, $A''=E^n/U$, of dimension $n/2$. Then $d_0=0$, $d_1=2n$, $d_2=n/2$, $d=5n/2$. The induced
map $\Psi''$ is onto $\C^2$, $V''=\Psi''^{-1}(V)$ is a hyperplane, and $W''=\ker\Psi''\subset V''$ is $\Qb$-rational of
dimension $d-2$. For a connected $G'\subset G''$ put $r'=\dim\Psi''(T_{G'})$; then $\tau=\delta+r'-2$, so only $r'\le1$
can satisfy $\delta>\tau$, and $\delta-\tau=2-r'$.

\emph{The data.} $w=(2\pi i,2ic\Omega)\in\Psi(\Omega_K(G))$. Suppose $x_1,x_2\in\mathcal S_K$ with $2\pi i,x_1,x_2$
$K$-independent, and put $v_j=(x_j,sx_j)\in V\cap\Psi(\Lambda_K(G))$. Let $Y''$ be generated by the images in $T_{G''}$
of elements of $\Lambda(G^n)$ with $\Phi$-images $D\theta_hw$ (taken in $\Omega(G^n)$, $D$ a common denominator),
$\theta_hv_1$, $\theta_hv_2$ ($h=1,\dots,n$). Then $\mathrm{rk}\,Y''=3n$, $\Gamma=\exp Y''\subset G''(\Qb)$, and
$\kappa=n$: a period in $Y''$ has first coordinate in $2\pi iK$, so its $v$-parts vanish. The inequality of
Theorem~\ref{thm:W88} reads $\lambda+\delta_1+2\delta_2\le(2-r')\cdot2n$.

\emph{No obstruction.}
\begin{itemize}
\item $G'=0$: $\lambda=2n$, so the left side is $2n+2n+n=5n>4n$.
\item $r'=1$ with image $\C(0,1)$ (this covers every $G'$ with an abelian part, since $\psi$ does not vanish on any
nonzero abelian subvariety of $A''$): an element of $Y''$ lying in $T_{G'}+\Omega(G'')$ has first coordinate in
$2\pi iK$, so $\mathrm{rk}(Y''\cap(T_{G'}+\Omega(G'')))\le n$ and $\lambda\ge2n$; the left side is $\ge3n>2n$.
\item $r'=1$ with image a line $\C(p,q)\neq\C(0,1)$: then $G'=T_L$ with $\dim L\le n$. Two $K$-independent data in
$T_{G'}+\Omega(G'')$ would put $x_1,x_2$ in a $K$-space of dimension $\le2$ containing $2\pi i$ (namely
$K\cdot2\pi i+K\pi^2/\Omega$ if $q=0$, and $K\cdot2\pi i+K\cdot p\,\Omega/(sp-1)$ otherwise, using $s\notin K$), against
the $K$-independence of $2\pi i,x_1,x_2$. So $\lambda\ge n$, $\delta_1\ge n$, $\delta_2=n/2$, and the left side is
$\ge3n>2n$.
\item $r'=0$ forces $G'=0$.
\end{itemize}
So no connected $G'\ne G''$ satisfies the conclusion of Theorem~\ref{thm:W88}, a contradiction.
\end{proof}

\noindent\emph{Controls.} If $E$ has no CM then $U=0$, $d_2=n$, and $G'=0$ satisfies the inequality ($6n\le6n$); so
does $G'=0$ with $\kappa$ deleted ($5n\le6n$), with only one datum, or with $k\not\subset K$. If $s\in K$ (a rational line) the
full subtorus over $V$ satisfies it with equality; replacing $E$ by $\Gm$ makes $s$ algebraic, and then the statement is
false ($s=2$, $x=\log2$).

\begin{corollary}[MG2: four exponentials at $t\in\Qb^\times\Omega/\pi$]\label{cor:MG2}
Let $t\in\Qb^\times\cdot\Omega/\pi$ and $\alpha,\beta\in\Qb^\times$ with $\log\alpha,\log\beta$ linearly independent over
$\Q$. Then $\alpha^t=e^{t\log\alpha}$ and $\beta^t=e^{t\log\beta}$ are not both algebraic. If moreover $\log\alpha,\log\beta,2\pi i$
are linearly independent over $\Q$, then $t\log\alpha$ and $t\log\beta$ are not both in $\Qb\Lg+\Qb\Lg_E$. For example,
$2^{\vp/\pi}$ and $3^{\vp/\pi}$ are not both algebraic.
\end{corollary}

\begin{proof}
If $\log\alpha,\log\beta,2\pi i$ are $\Q$-independent, take $x_1=\log\alpha$, $x_2=\log\beta$ in Theorem~\ref{thm:MG}:
with $2\pi i$ they are three $\Q$-independent logarithms, hence $K$-independent. Otherwise
$m\log\beta=a\log\alpha+2\pi i\,b$ with integers $m\ne0$, $b\ne0$. If $\alpha^t$ and $\beta^t$ are algebraic, so is
$e^{2\pi ibt}=(\beta^t)^m(\alpha^t)^{-a}$, i.e.\ $2\pi ibt=2ib\,c\,\Omega\in\Lg$ (with $t=c\Omega/\pi$). Then
$\Omega\in\Qb\Lg$, which the analytic subgroup theorem \cite{Wu89} excludes (on $\Gm\times E$, the smallest algebraic
subgroup whose tangent space contains $(\ell,\Omega)$ is $\Gm\times E$). (Remark~\ref{rem:MG5d} gives a proof that needs
no such step.)
\end{proof}

\noindent The determinations matter: $\log\alpha=\log2$ and $\log\beta=\log2+2\pi i$ are $\Q$-independent, and for the
determinations $e^{t\log2}$ of both powers the statement would say $2^t\notin\Qb$, a one-exponential statement that we do
not prove.

\begin{remark}[Second derivation via \S5d]\label{rem:MG5d}
Theorem~\ref{thm:W88k} gives Theorem~\ref{thm:MG} with a shorter subgroup analysis, and the first sentence of
Corollary~\ref{cor:MG2} without the case distinction. Keep $G''=\Gm^{2n}\times A''$, $\Psi''$, $V''=\Psi''^{-1}(V)$ and
$W''=\ker\Psi''$ from the proof of Theorem~\ref{thm:MG}: $d=5n/2$, $d_0=0$, $d_1=2n$, $d_2=n/2$, and $W''$ is
$\Qb$-rational of dimension $t=d-2$, so $d-1-t=1$. Let $Y''$ be generated by the $3n$ lifts of $D\theta_hw$,
$\theta_hv_1$, $\theta_hv_2$ ($h=1,\dots,n$). The $n$ period lifts lie in $Y''\cap\ker\exp$, so $\kappa\ge n$ (in fact
$\kappa=n$), and (5.4) reads $3n>(2n+n-\kappa)\cdot1$, i.e.\ $3n>2n$.

\emph{The rank $m=3n$.} For the data of Theorem~\ref{thm:MG} ($2\pi i,x_1,x_2$ $K$-independent) the first coordinates of
the $\Psi''$-images $D\theta_h\cdot2\pi i$, $\theta_hx_j$ are $\Q$-independent. For the data of Corollary~\ref{cor:MG2}
($x_j=\log\alpha_j$, $sx_j=t\log\alpha_j\in\Lg$) only the $\Q$-independence of $\log\alpha,\log\beta$ is needed. Take the
period lifts with toric part $(2\pi iD e_h,0)$ and $E^n$-part $m^{(h)}\Omega\in\Lambda_E^n$, where
$\psi(m^{(h)})=2icD\theta_h$ (possible by Lemma~\ref{lem:K}, as $ic\in K$; $D$ a common denominator), and the data lifts toric, $(x_je_h,\ sx_je_h)$. The linear form induced
by $\theta^{\mathsf T}z$ on $T_{A''}$ takes the $\Q$-independent values $2icD\theta_h\Omega$ on the period lifts and $0$ on the
data lifts, so a $\Q$-relation among the $3n$ lifts involves only the data lifts; their first-block coordinates $x_je_h$
are $\Q$-independent.

\emph{No subgroup in $V''$.} Theorem~\ref{thm:W88k} gives a connected $G'\ne0$ with $T_{G'}\subset V''$, and $G'=T_L\times
B_S$. If $B_S\ne0$, then $T_{G'}$ contains a vector $(0,b)$ with $\theta^{\mathsf T}b\ne0$ (since $\psi$ does not vanish on any
nonzero abelian subvariety of $A''$), whose image $(0,\theta^{\mathsf T}b)$ is not on $V=\C(1,s)$. If $B_S=0$, then $L\ne0$, and
for $0\ne l=(l_1,l_2)\in L$ the image $(\theta^{\mathsf T}l_1,\theta^{\mathsf T}l_2)\in K^2\setminus0$ lies on $\C(1,s)$; so
$\theta^{\mathsf T}l_1\ne0$ and $s\in K$, which is false ($\Omega/\pi\notin\Qb$ \cite{Ma75}).

For $K=\Q(i)$, $n=2$, $\theta=(1,i)$: $A''\cong E$ via $(z_1,z_2)\mapsto z_1+iz_2$, so $G''=\Gm^4\times E$, $d=5$, $t=3$,
$m=6$, $\kappa=2$, and (5.4) is $6>6-\kappa$. This route uses neither the $K$-independence of $2\pi i,\log\alpha,\log\beta$
nor W\"ustholz's theorem: it proves the first sentence of Corollary~\ref{cor:MG2} directly, also when $2\pi i$ lies in the
$\Q$-span of $\log\alpha,\log\beta$. Its subgroup exclusion uses only $s\notin K$. It rests on Theorem~\ref{thm:W88k} as
published, where the main proof rests on Theorem~\ref{thm:W88} and Lemma~\ref{lem:U}.
\end{remark}

\begin{corollary}[MG1]\label{cor:MG1}
For every CM curve $E$, every period $\Omega$ and all $a,b,c\in\Qb^\times$, at least one of
$e^{a\pi^2}$, $e^{b\pi\Omega}$, $e^{c\pi^2/\Omega}$ is transcendental. On $E_i$: one of $e^{\pi^2}$, $e^{\pi\vp}$,
$e^{\pi\eta(\vp)}$ is transcendental.
\end{corollary}

\begin{proof}
Take $s=\Omega/\pi$ (up to $\Qb^\times$), $x_1=\pi^2$, $x_2=\pi^2/\Omega$: if all three exponentials are algebraic then
$x_1,x_2\in\Qb\Lg$ and $sx_1=\pi\Omega$, $sx_2=\pi$ lie in $\Qb\Lg$. The numbers $2\pi i,\pi^2,\pi^2/\Omega$ are
$\Qb$-independent. A relation without the $\pi^2/\Omega$ term contradicts the transcendence of $\pi$. Otherwise
$\pi/\Omega=a_1+b_1\pi$ with $a_1,b_1\in\Qb$. If $b_1=0$ then $\Omega/\pi\in\Qb$, which is false. If $b_1\ne0$ then
$\pi\Omega=(\pi-a_1\Omega)/b_1$; since $\pi\Omega\in\Qb\Lg$ (as $e^{b\pi\Omega}\in\Qb$) and $\pi\in\Qb\Lg$, this gives
$a_1\Omega\in\Qb\Lg$, so $a_1=0$ because $\Omega\notin\Qb\Lg$ (proof of Corollary~\ref{cor:MG2}); then $\Omega=1/b_1\in\Qb$,
which is false. (Chudnovsky's algebraic independence of $\pi$ and $\Omega$ \cite{Ch84} would give this at once, but is not
needed.) On $E_i$, $\pi/\vp=\eta(\vp)$.
\end{proof}

\noindent For a general CM curve, $\eta(\Omega)=\kappa_1\Omega+\kappa_2\pi/\Omega$ with $\kappa_1,\kappa_2\in\Qb$ and $\kappa_2\ne0$ (Legendre);
$\kappa_1=0$ when $j(E)\in\{0,1728\}$. Since $\pi\Omega,\pi\eta\in\Qb\Lg$ imply $\pi^2/\Omega=(\pi\eta-\kappa_1\pi\Omega)/\kappa_2\in\Qb\Lg$,
the statement ``one of $e^{\pi^2}$, $e^{\pi\Omega}$, $e^{\pi\eta}$ is transcendental'' holds for every CM period as well.

\begin{corollary}[MG3]\label{cor:MG3}
If $e^{\pi^2}$ and $e^{c\pi\Omega}$ ($c\in\Qb^\times$) are algebraic, then $\alpha^{\Omega/\pi}\notin\Qb$, for every determination,
for every $\alpha\in\Qb^\times$ such that $2\pi i,\pi^2,\log\alpha$ are linearly independent over $\Q$.
\end{corollary}

\begin{proof}
With $s=\Omega/\pi$, take $x_1=\pi^2$ (then $sx_1=\pi\Omega\in\Qb\Lg$) and $x_2=\log\alpha$ (then $sx_2\in\Lg$ if
$\alpha^{\Omega/\pi}\in\Qb$). These are $\Q$-independent logarithms together with $2\pi i$, hence $K$-independent, and
Theorem~\ref{thm:MG} gives a contradiction. The hypothesis is needed: for $c=1$ and $\alpha=e^{\pi^2}$ ($\log\alpha=\pi^2$) one has
$\alpha^{\Omega/\pi}=e^{\pi\Omega}\in\Qb$. (Corollary~\ref{cor:MG3} is also the case $\beta=e^{\pi^2}$ of the second sentence of
Corollary~\ref{cor:MG2}.)
\end{proof}

\noindent\emph{Overlap with known results.} Waldschmidt's \cite[Cor.~2.7]{W07} (after Diaz \cite{Di04}) says: if
$x_1,x_2\in\R\cup i\R$ are linearly independent over $\Q$ and $y_1,y_2,\bar y_2$ are linearly independent over $\Q$, then one of the four numbers
$e^{x_iy_j}$ is transcendental. With $x=(1,t)$ it gives Corollaries~\ref{cor:MG2} and~\ref{cor:MGn2} for $t\in\R\cup
i\R$ and pairs with $\log\alpha,\log\beta,\overline{\log\beta}$ $\Q$-independent (e.g.\ $(2,1+2i)$). With the roles
exchanged, $x=(\log\alpha,\log\beta)$ and $y=(1,t,\bar t)$, it gives them for pairs with logarithms in $\R\cup i\R$
(e.g.\ $(2,3)$) whenever $1,t,\bar t$ are $\Q$-independent, e.g.\ $t=e^{i\pi/5}\vp/\pi$. Rescaling widens the
first orientation to every $t$ in our classes. Such a $t$ is $\gamma r$ with $\gamma\in\Qb^\times$ and $r\in\R$
($\bar E$ has CM by the same field, so $\bar\Omega\in\Qb^\times\Omega$; and $\pi$ is real). If $\alpha^t,\beta^t\in\Qb$,
then $y=(\log\alpha,\log\beta)$ satisfies $y_j,\ ry_j=\gamma^{-1}ty_j\in\Lt$, and the $\Lt$-version \cite[Thm.~4.1]{W07} at the real
point $r$ excludes every pair with $\log\alpha,\log\beta,\overline{\log\beta}$ linearly independent (over $\Qb$, which
for logarithms is the same as over $\Q$ by Baker's theorem), whatever the argument of $t$. The remaining pairs have a
conjugation-stable $\Q$-span of logarithms, which has a basis in $\R\cup i\R$. For them the second orientation applies as
soon as $t\notin\R\cup i\R$: then $1,t,\bar t$ are $\Q$-independent, because $r$ is transcendental and $\bar\gamma/\gamma\notin\Q$.
So at $t\notin\R\cup i\R$ Corollaries~\ref{cor:MG2} and~\ref{cor:MGn2} are already known. What remains new is real or
purely imaginary $t$ (as for $\vp$, $\vp/\pi$, $\Omega_\rho$) with a conjugation-stable span of logarithms (such as
$(2,3)$, or $(5,(1+2i)/(1-2i))$); Corollary~\ref{cor:MG1}; and the dimension statements themselves. We found no case of the
four exponentials conjecture at a fixed transcendental $t$ of this kind in \cite{W22} or \cite[p.~27]{SMF84}. The
closest published theorem is Waldschmidt's \cite[Thm.~4.2]{W05}, a mixed statement on $\Ga\times\Gm\times E$ without CM
whose second row consists of elliptic logarithms; it needs seven columns (one of six numbers $\ell_1u_j-\ell_ju_1$ is
transcendental) in general, and five columns (one of four) in its period-column case. CM and the mixing map bring the count down to two
data with ordinary logarithms.

\subsection{The fibre twin}
Let $E^\natural$ be read on its conjugate-eigen fibre coordinate $w$ (\S\ref{sec:MX}), with $\eta^*=\eta^*(\Omega)$. By
Legendre and $\eta^*(a\lambda)=\bar a\eta^*(\lambda)$, $\Omega\eta^*\in\Qb^\times\pi$ for every CM curve. Call
\emph{fibre values} the $w$-coordinates of $\Lambda(E^\natural)$ (that is, $\zeta^*(u)+\Qb$ for $u\in\Lg_E\setminus\Lambda_E$,
$\eta^*(\lambda)+\Qb$ at lattice points).

\begin{theorem}[Theorem MG$^\natural$]\label{thm:MGn}
Let $c\in\Qb^\times$ and $s=c\,\eta^*/\pi\in\Qb^\times/\Omega$. For every number field $K\supseteq k$,
$\dim_K\{x\in K\Lg:\ s\,x\in K\Lg+K\cdot(\text{fibre values})\}\le2$.
\end{theorem}

\begin{proof}
As in the proof of Theorem~\ref{thm:MG} we may enlarge $K$; we take $K\supseteq k(i,c)$ with $n=[K:\Q]\ge4$ (for $n=2$
the count below does not close). Let $G=\Gm\times\Gm\times E^\natural$, $\Psi(x,y,X)=(x,\,y+w(X))$, $V=\C(1,s)$, and
$\theta$, $p_\theta$, $\Phi=\Psi\circ p_\theta$ as before.

\emph{The group.} There are no nonzero homomorphisms between tori and connected subgroups of $(E^\natural)^n$ in either
direction, so connected subgroups of $\Gm^{2n}\times(E^\natural)^n$ are products. A connected $H\subset(E^\natural)^n$ over
$B_S$ contains the canonical lift $B_S^\natural$ (since $\mathrm{Hom}(B_S^\natural,\Ga)=0$), whose tangent space is
$\{(\sigma(M)z,\bar\sigma(M)w)\}$; so $\theta^{\mathsf T}w$ vanishes on $T_H$ only if $\phi(S)=0$. Hence the largest
connected $U$ with $T_U\subset\ker\Phi$ is $U=B^\natural_{\ker\phi}+\{\theta^{\mathsf T}w=0\}$, and $G''=\Gm^{2n}\times Q''$ with
$Q''=(E^\natural)^n/U$, an extension of $A''=E^n/B_{\ker\phi}$ (dimension $n/2$) by $\Ga$ (coordinate
$\sigma=\theta^{\mathsf T}w$). Over an abelian subvariety $B'\ne0$ of $A''$, coming from $S'\supsetneq\ker\phi$, we have
$\phi(S')\neq0$, so the image of $B_{S'}^\natural$ contains the fibre. Thus $Q''$ has no splitting over any $B'\ne0$, and its
connected subgroups are $0$, the fibre $F$, and the $\pi^{-1}(B')$ (dimension $\le n/2+1$). $Q''$ is a quotient of an
anti-affine group, so $\mathrm{Hom}(Q'',\Ga)=0$: $d_0=0$, $d_1=2n$, $G_2=Q''$, $d_2=n/2+1$. The induced map is
$\Psi''(t_1,t_2,q)=(\theta^{\mathsf T}t_1,\theta^{\mathsf T}t_2+\sigma(q))$, $V''$ is a hyperplane, $W''=\ker\Psi''$, and
$\tau=\delta+r'-2$ as before.

\emph{The data.} The period datum is $(2\pi i,2ic\eta^*)\in\Psi(\Omega_K(G))$ (using $ic\in K$ and that $\phi$ is onto
$K$). With $x_1,x_2$ as in Theorem~\ref{thm:MG}, $\mathrm{rk}\,Y''=3n$ and $\kappa=n$ (a period of $G''$ has first coordinate
in $2\pi iK$). The inequality reads $\lambda+\delta_1+2\delta_2\le(2-r')(2n+2)$.

\emph{No obstruction.} Every $G'$ is $T_L\times H$ with $H\in\{0,F,\pi^{-1}(B')\}$, and $H\ne0$ forces
$\C(0,1)\subset\Psi''(T_{G'})$.
\begin{itemize}
\item $G'=0$: $\lambda=2n$, and the left side is $2n+2n+(n+2)=5n+2>4n+4$.
\item Image $\C(0,1)$ ($H\ne0$, or $L$ in the second block; $l=\dim L\le n$, $h=\dim H\le n/2+1$): elements of $Y''$ in
$T_{G'}+\Omega(G'')$ have first coordinate in $2\pi iK$, so $\lambda\ge2n$, and the left side is
$\ge2n+(2n-l)+2(n/2+1-h)\ge3n>2n+2$.
\item $H=0$ and image a line $\C(1,q)$, $q\in K$: a datum $X$ in $T_{G'}+\Omega(G'')$ satisfies
$X(s-q)\in K\cdot2\pi i+K\cdot2\pi is$, so $X\in K\cdot2\pi i+K\cdot2\pi i/(s-q)$ (here $s\notin K$ because $\Omega$ is
transcendental). So $\lambda\ge n$, and the left side is $\ge n+n+(n+2)=3n+2>2n+2$.
\item $r'=0$ forces $G'=0$; $r'=2$ fails $\delta>\tau$.
\end{itemize}
The minimum margin is $n-2$.
\end{proof}

\begin{corollary}[MG$^\natural$2: four exponentials at $t\in\Qb^\times\Omega$]\label{cor:MGn2}
If $\log\alpha,\log\beta$ are linearly independent over $\Q$ and $c\in\Qb^\times$, then $\alpha^{c\Omega}=e^{c\Omega\log\alpha}$
and $\beta^{c\Omega}=e^{c\Omega\log\beta}$ are not both algebraic. For example, $2^{\vp}$ and $3^{\vp}$ (real logarithms) are
not both algebraic.
\end{corollary}

\begin{proof}
Put $x_j=c\Omega\log\alpha_j$ (with $\alpha_1=\alpha$, $\alpha_2=\beta$). If $\alpha^{c\Omega},\beta^{c\Omega}$ are
algebraic then $x_j\in\Lg$, and $sx_j\in\Qb\log\alpha_j\subset\Qb\Lg$. A $K$-relation among $2\pi i,x_1,x_2$ without the
$2\pi i$ term contradicts Baker; one with it gives $0\ne\ell\in\Qb\Lg\cap\Qb^\times\eta^*$, impossible by the analytic
subgroup theorem on $\Gm^r\times E^\natural$ \cite{Wu89}. (Remark~\ref{rem:MGn5d} gives a proof that needs no such step.)
\end{proof}

\begin{remark}[Second derivation via \S5d]\label{rem:MGn5d}
As in Remark~\ref{rem:MG5d}, keep $G''=\Gm^{2n}\times Q''$, $\Psi''$, $V''$ and $W''=\ker\Psi''$ from the proof of
Theorem~\ref{thm:MGn}: $d=5n/2+1$, $d_0=0$, $d_1=2n$, $d_2=n/2+1$, $t=d-2$. With the $3n$ lifts of $D\theta_hw$,
$\theta_hv_1$, $\theta_hv_2$ one has $m=3n$ and $\kappa\ge n$ (in fact $\kappa=n$), and (5.4) reads
$3n>(2n+n+2-\kappa)\cdot1$, i.e.\ $3n>2n+2$. It holds as soon as $\kappa\ge3$, so for every admissible $n\ge4$ (at $n=2$ it
would read $6>6$). For $K=\Q(i,\sqrt2)$, $n=4$: $G''=\Gm^8\times Q''$, $d=11$, $t=9$, $m=12$, $\kappa=4$, and (5.4) is
$12>14-\kappa$.

\emph{The rank.} For the data of Theorem~\ref{thm:MGn} it follows from the first coordinates, as in
Remark~\ref{rem:MG5d}. For Corollary~\ref{cor:MGn2} take $s=1/(c\Omega)\in\Qb^\times\eta^*/\pi$, so that
$sx_j=\log\alpha_j\in\Lg$, and take the data lifts toric, $(x_je_h,\ \log\alpha_j\,e_h)$. Take the period lifts with toric
part $(2\pi iDe_h,0)$; the second coordinate $D\theta_h\cdot2\pi is$ of their image is then carried by the fibre
coordinate $\sigma$ of $Q''$. So $\sigma$ separates the period lifts from the data lifts, and the first-block coordinates
$x_je_h$ of the data lifts are $\Q$-independent as soon as $\log\alpha,\log\beta$ are.

\emph{No subgroup in $V''$.} A connected $G'\ne0$ with $T_{G'}\subset V''$ is $T_L\times H$ with $H\in\{0,F,\pi^{-1}(B')\}$. If
$H\ne0$, then $F\subset H$, and $T_F$ contains a vector with image $(0,1)\notin V$. If $H=0$, then $L\ne0$ and, as in
Remark~\ref{rem:MG5d}, a nonzero $l\in L$ gives $s\in K$, which is false since $s\in\Qb^\times/\Omega$ and $\Omega$ is
transcendental.

So this route, too, needs only the $\Q$-independence of $\log\alpha,\log\beta$ and no W\"ustholz step. It rests on
Theorem~\ref{thm:W88k} as published.
\end{remark}

\noindent By the symmetry $t\leftrightarrow1/t$, Corollaries~\ref{cor:MG2} and~\ref{cor:MGn2} give the four exponentials
conjecture at every $t\in\Qb^\times\Omega/\pi\cup\Qb^\times\pi/\Omega$ and at every $t\in\Qb^\times\Omega\cup\Qb^\times/\Omega$,
for every
CM period $\Omega$.

\subsection{The affine version and strong four exponentials}\label{sec:MGaff}
Adding one factor $\Ga$ to each coordinate of $\Psi$ (so $G=\Ga^2\times\Gm^2\times E$ with coordinates $(a,a',x,y,z)$ and
$\Psi(a,a',x,y,z)=(x+a,\,y+z+a')$) gives the following statement. In the quotient
$G''=\Ga^2\times\Gm^{2n}\times A''$ one has $d_0=2$, which does not enter the left-hand side of Theorem~\ref{thm:W88}. The
only new subgroups are $\Ga$-lines $\{(\mu a,a)\}$ with $\mu\notin K$, which are excluded with margin $n$. So the count is
the same as for Theorem~\ref{thm:MG}; Lemma~\ref{lem:U} applies since $5n+2\ge3n+4$. This was checked at proof-sketch level
by both internal referees of Theorem~\ref{thm:MG} and by a third; the subgroup list is not written out here. The base change
in $K$ now uses Baker's inhomogeneous theorem.

\begin{theorem}[Affine MG; proof sketch]\label{thm:MGaff}
With $E$, $\Omega$, $c$, $s=c\Omega/\pi$ and $K$ as in Theorem~\ref{thm:MG},
\[
\dim_K\{x\in K+K\Lg:\ s\,x\in K+K\Lg+K\Lg_E\}\le2 .
\]
\end{theorem}

\begin{corollary}[proof sketch]\label{cor:MGaff}
\begin{itemize}
\item[(i)] For $a,b,c\in\Qb^\times$, one of $e^{a\pi^2}$, $e^{b\pi\Omega}$, $e^{c\Omega/\pi}$ is transcendental.
\item[(ii)] For $\alpha\in\Qb^\times$ not a root of unity, $e^{\Omega/\pi}$ and $\alpha^{\Omega/\pi}$ (any determination) are not
both algebraic.
\item[(iii)] (Strong four exponentials at ratio $t\in\Qb^\times\Omega/\pi$.) If $y_1,y_2$ are linearly independent over
$\Qb$, then $y_1$, $y_2$, $ty_1$, $ty_2$ are not all in $\Lt=\Qb+\Qb\Lg$. By $y\mapsto ty$ the same holds at
$t\in\Qb^\times\pi/\Omega$. (The strong statement is invariant under $t\mapsto ct$, $c\in\Qb^\times$, and $\Qb^\times\Omega/\pi$ contains
a real number. So this is a statement at real $t$, where the $\Lt$-version of the conjugation argument
\cite[Thm.~4.1]{W07} applies when $y_1,y_2,\bar y_2$ are $\Qb$-independent. The new content is the remaining $y$.)
\end{itemize}
\end{corollary}

\begin{proof}[Sketch]
(i) Take $x_1=1$ ($s\cdot1\in\Qb\Lg$ if $e^{c\Omega/\pi}$ is algebraic, with $s$ adjusted by a constant) and $x_2=\pi^2$. With $2\pi i$
these are $K$-independent, because $1,\pi,\pi^2$ are $\Qb$-independent (Lindemann). (ii) Take $x_1=1$ and $x_2=\log\alpha$;
then $1,2\pi i,\log\alpha$ are $K$-independent by Baker's theorem. (iii) Take $K$ containing all the coefficients. Then
$y_1,y_2\in T_K:=\{x\in K+K\Lg:\ sx\in K+K\Lg+K\Lg_E\}$, which has dimension $\le2$ and contains $2\pi i$. Since $y_1,y_2$ are
$K$-independent, $2\pi i\in Ky_1+Ky_2$. Then $2ic\Omega=s\cdot2\pi i\in\Lt$, which contradicts the analytic subgroup theorem
($\Omega\notin\Qb+\Qb\Lg$).
\end{proof}

\subsection{The general form, and elliptic four exponentials}\label{sec:MGgen}
Theorems~\ref{thm:MG} and~\ref{thm:MGn} are the cases $\sigma_1=\sigma_3=\Gm$ of the following count. Let
$\sigma_1,\sigma_2,\sigma_3$ be ``slots'': $\Gm$ (period $2\pi i$), the base of a CM curve (period $\Omega$), or the fibre of
the universal vector extension of a CM curve (period $\eta^*$). In the count of Theorem~\ref{thm:W88} the first two have
weight $1$ per factor and a fibre slot has weight $1+2/n$. Write $\Lambda_K(\sigma)$ for the
$K$-span of the corresponding logarithms ($\Lg$, $\Lg_E$, or fibre values). Let $p_j$ be the period of $\sigma_j$,
$r=p_2/p_1$, and take $\Psi$ reading $\sigma_1$ into the first coordinate and $\sigma_2+\sigma_3$ into the second.

\begin{theorem}[General form; proof sketch]\label{thm:MGgen}
Assume:
\begin{itemize}
\item $\sigma_1,\sigma_2$ are not of the same kind on isogenous groups (so $r\notin\Qb$);
\item if $\sigma_2$ and $\sigma_3$ are of the same kind on isogenous groups, they are merged into one slot;
\item (allowing $r\in\Qb^\times p_2/p_1$) $K$ contains this constant, the CM fields, and the Legendre and isogeny constants,
and $[K:\Q]\ge2f+2$, where $f$ is the number of fibre slots;
\item $\sigma_1,\sigma_3$ are not both fibre slots.
\end{itemize}
Then
\[
\dim_K\{\lambda\in\Lambda_K(\sigma_1):\ r\lambda\in\Lambda_K(\sigma_2)+\Lambda_K(\sigma_3)\}\le2 .
\]
\end{theorem}

\noindent The count is that of Theorem~\ref{thm:MG}, with the threshold $n>2f$ coming from $G'=0$ ($5n+2f>4n+4f$). An
internal referee checked all $27$ slot-kind triples by an exact integer case check. The load-bearing step (a bound on the
data in subtori with image a $K$-line) was done by hand, and the subgroup lists are not written out in this paper. The
excluded case is ``four exponentials for quasi-logarithms'': $\sigma_1,\sigma_3$ fibre slots of isogenous curves (for
non-isogenous curves we have not written it out). It has mixed subgroups whose fibre slope $\mu$ is algebraic but possibly
outside $K$, and it needs the $K(\mu)$-linear independence of $K$-independent fibre values, which we have not proved. When $\sigma_1=\sigma_3$ is the base of a CM
curve $E$, Theorem~\ref{thm:MGgen} is an elliptic analogue of four exponentials:

\begin{corollary}[Elliptic four exponentials at CM slopes; proof sketch]\label{cor:ell4E}
Let $E$ have CM, with period $\Omega$, and let $t$ be one of $\pi/\Omega^2$, $\Omega'/\Omega$, $\eta^*(\Omega')/\Omega$, where $\Omega'$ is
a period of a CM curve $E'$ not isogenous to $E$, all up to $\Qb^\times$. If $u_1,u_2\in\Lg_E$ and $\Omega$ are linearly
independent over $k$, then $tu_1$ and $tu_2$ do not both lie in $\Qb\Lg_E+\Qb\Lambda(\sigma_2)$. Here $\Lambda(\sigma_2)$ is the
set of fibre values of $E$ (for $t=\pi/\Omega^2$), $\Lg_{E'}$ (for $t=\Omega'/\Omega$), or the set of fibre values of $E'$ (for
$t=\eta^*(\Omega')/\Omega$). In particular $tu_1,tu_2$ are not both elliptic logarithms of algebraic points of $E$.
\end{corollary}

\section{Theorem M}\label{sec:M}

\subsection{First proof (the count)}
Let $G=E\times\Gm\times\Ga$, $\Psi(e,t,a)=(e,t+a)$, $V=\C(\Omega,i\pi)$, $W=0$, and $K\supseteq k$ containing the
algebraic coefficients involved. Then $\alpha_K(G)=1+1+0=2$; $\Psi(\Omega_K(G))=K\Omega\times K\cdot2\pi i$ since
$K\Lambda_E=K\Omega$ for CM; so $\kappa=1$. Condition (i): $V$ has a nonzero rational point iff $\Omega/\pi\in\Qb$, which is
false (Masser \cite{Ma75}; cf.\ \cite[Thm.~26]{W06}; also Schneider). Theorem~\ref{thm:RW} gives
$\dim_K(V\cap\Psi(\Lambda_K))\le1$, and the period $(\Omega,i\pi)$ lies in it. If $i\pi u=\Omega\ell$ then
$(u,\ell)\in V\cap\Psi(\Lambda_K)$, so $(u,\ell)=s(\Omega,i\pi)$ with $s\in K$. If $u\in\Lg_E$, Masser's theorem
(\cite[Thm.~27]{W06}: elements of $\Lg_E$ that are $k$-independent are $\Qb$-independent) gives $s\in k$. \qed

\smallskip\noindent\emph{Controls.} If $E$ has no CM, $\alpha_K(E)=2$ and the count ties (no conclusion). Replacing $E$ by
$\Gm$ ($\Omega\to2\pi i$) makes $V$ rational, so (i) fails --- as it must, since $u=2\ell$ solves $i\pi u=2\pi i\ell$. For
$e^{\pi^2}$ itself ($G=\Ga\times\Gm^2$, $V=\C(1,i\pi)$) one finds $\kappa=0$ and an exact tie.

\subsection{Second proof for $u\in\Lg_E$ (Laurent's identity and W\"ustholz)}\label{sec:M2}
We treat $u\in\Lg_E$, which covers every item of Corollary~\ref{cor:M}; a version for $u\in\Qb\Lg_E$ would need several
points $P_j$ and is not written out here.
Let $\Lambda_E=\Omega\,\Z[\alpha]$ (a sublattice of finite index suffices), and write Masser's relation \cite{Ma75} as
$\eta(\alpha\Omega)=\bar\alpha\eta(\Omega)+c\,\Omega$ with $c\in\Qb$. Legendre's relation gives
$(\alpha-\bar\alpha)\eta(\Omega)-c\Omega=2\pi i/\Omega$, and $f(z)=\zeta(\alpha z)-\bar\alpha\zeta(z)-cz$ is an elliptic
function defined over $\Qb$. With the third-kind period $\lambda(u,\Omega)=\Omega\zeta(u)-\eta(\Omega)u$ one obtains, for
$u\in\Lg_E$ not a pole of $f$,
\begin{equation}\label{eq:laurent}
\frac{2\pi i\,u}{\Omega}=\bar\alpha\,\lambda(u,\Omega)-\lambda(\alpha u,\Omega)+\Omega f(u),\qquad f(u)\in\Qb .
\end{equation}
On $E_i$ this reads $2\eta(\vp)u=i\lambda(iu,\vp)-\lambda(u,\vp)$ (cf.\ \cite[\S9, p.~160]{La82}, where the general CM
form is derived; the display above is its case $\mu=i$). Hence $i\pi u=\Omega\ell$ with
$\ell\in\Lt$ is a $\Qb$-linear relation among $1$, logarithms of algebraic numbers, $\Omega$, $\lambda(u,\Omega)$ and
$\lambda(\alpha u,\Omega)$. These are coordinates of a point of $\Lambda(\mathcal G)$ for
$\mathcal G=G_P\times G_{\alpha P}\times\Gm^r\times\Ga$, where $G_P$ is the extension of $E$ by $\Gm$ attached to
$P=\exp_E(u)$ (Serre \cite{Se79}) --- add a factor $E^\natural$ when algebraic constants enter. By the analytic subgroup theorem \cite{Wu89}, the
smallest algebraic subgroup $B$ whose tangent space contains this point lies in the hyperplane defined by the relation.
$B$ surjects onto the diagonal of $E\times E$, and over it the toric part carries the extension classes
$[a+b\alpha]P$ for characters $(a,b)$; these vanish only for $a=b=0$ when $P$ is not torsion. So $B$ contains both fibre
tori, on which the relation has nonzero coefficients --- a contradiction. Hence $P$ is torsion, which is the trivial
case. \qed

\subsection{The corollaries}
(M1): $u=c\vp/\pi$, $\ell=ic$ gives $s=c/\pi\notin\Qb$. (M2): $\ell=i\log\alpha$ gives $\log\alpha\in\Qb\pi$, so
$\alpha=(-1)^{-is}$ with $s$ algebraic, a root of unity by Gel'fond--Schneider. (M3): if $e^{\pi^2}\in\Qb$ then
$\ell=ic\pi^2\in\Lt$ and $u=c\pi\Omega$ give $s=c\pi\notin\Qb$. (M3$'$): $u=\vp\log\alpha$, $\ell=i\pi\log\alpha$ give
$\log\alpha\in\Qb^\times$, contradicting Hermite--Lindemann; (M3) is the case $\alpha=-1$, using $\wp(iz)=-\wp(z)$. (M$\eta$):
$\eta(\vp)=\pi/\vp$. On a general CM curve $\eta(\Omega)=\kappa_1\Omega+\kappa_2\pi/\Omega$ with $\kappa_2\ne0$, and the
statement becomes $(\eta-\kappa_1\Omega)u\notin\Lt$.

\section{Theorem MX}\label{sec:MX}

\subsection{Set-up}
For a CM curve $E$ let $E^\natural$ be its universal vector extension, with base coordinate $z$ and conjugate-eigen
fibre coordinate $w$ as in Lemma~\ref{lem:natural}. In these coordinates the periods are $(\lambda,\eta^*(\lambda))$,
$\lambda\in\Lambda_E$, where $\eta^*(\lambda)=\eta(\lambda)-\kappa_1\lambda$ satisfies $\eta^*(a\lambda)=\bar a\,\eta^*(\lambda)$;
the logarithms of points of $E^\natural(\Qb)$ are $(u,\zeta^*(u)+t)$ with $u\in\Lg_E$, $t\in\Qb$ and
$\zeta^*=\zeta-\kappa_1\,\mathrm{id}$. On $E_i$ and $E_\rho$, $\kappa_1=0$. Since the characters $a\mapsto a,\bar a$ of $k$
are $K$-independent for $K\supseteq k$,
\begin{equation}\label{eq:split}
\Omega_K(E^\natural)=K\lambda_0\times K\eta^*(\lambda_0)\qquad(\Lambda_E\otimes\Q=k\lambda_0).
\end{equation}

\subsection{Proof of Theorem~\ref{thm:MX}}
Let $G=E_i^\natural\times E_\rho^\natural\times\Ga^2$, let $K\supseteq\Q(i,\sqrt{-3},\beta)$ be a number field, $n=[K:\Q]$
(so $4\mid n$), and let
\[
\Psi(x_i,x_\rho,a_1,a_2)=\bigl(z(x_i)+z(x_\rho)+a_1,\;w(x_i)+w(x_\rho)+a_2\bigr),\qquad V=\C\,w_1,
\]
with $w_1=(\vp,\beta\eta_\rho)$, $\eta_i=\eta(\vp)=\pi/\vp$, $\eta_\rho=\eta(\Omega_\rho)=2\pi/(\sqrt3\,\Omega_\rho)$, and
$w_2=(\Omega_\rho,\beta\tfrac{2}{\sqrt3}\eta_i)$. By (\ref{eq:split}), $w_1,w_2\in\Psi(\Omega_K(G))$ and
\[
\Psi(\Omega_K(G))=P:=(K\vp+K\Omega_\rho)\times(K\eta_i+K\eta_\rho);
\]
$w_1\parallel w_2$ because $\vp\cdot\tfrac{2}{\sqrt3}\eta_i=\tfrac{2\pi}{\sqrt3}=\Omega_\rho\eta_\rho$, and $w_1,w_2$ are $K$-independent
since $\Omega_\rho/\vp\notin\Qb$ (periods of non-isogenous curves are $\Qb$-independent \cite{Wu89}).

\begin{proposition}\label{prop:D}
$V\cap\Psi(\Lambda_K(G))=Kw_1+Kw_2$.
\end{proposition}

Granting Proposition~\ref{prop:D}: the point $(X,Y,a,t)$ lies in $\Lambda(G)$ and the hypothesis says that its image
lies on $V$, so $z(X)+z(Y)+a=x\vp+x'\Omega_\rho$ with $x,x'\in K$. By the linear independence of $1$ and elliptic
logarithms on non-isogenous curves (W\"ustholz \cite{Wu89}), $a=0$, $z(X)=x\vp$, $z(Y)=x'\Omega_\rho$, and by
Masser's theorem \cite{Ma75} $x\in\Q(i)$, $x'\in\Q(\rho)$. \qed

\smallskip
\noindent\emph{Proof of Proposition~\ref{prop:D}.} Suppose $v\in V\cap\Psi(\Lambda_K(G))$ with $v\notin Kw_1+Kw_2$.

\emph{Step 1: an elementary lemma.} For $L\in\{0,\ \C(1,0),\ \C(0,1)\}$ we have $V\cap(L+P)=Kw_1+Kw_2$. Indeed, if
$sw_1\in L+P$ with $L\ne\C(1,0)$ then $s\vp\in K\vp+K\Omega_\rho$, so $s\in K+K\Omega_\rho/\vp$ and $(\Omega_\rho/\vp)w_1=w_2$;
if $L=\C(1,0)$ then $s\beta\eta_\rho\in K\eta_i+K\eta_\rho$, so $s\in K+K\eta_i/\eta_\rho$ and
$(\eta_i/\eta_\rho)w_1=(\pi/\eta_\rho,\beta\eta_i)=\tfrac{\sqrt3}{2}w_2$.

\emph{Step 2: the group.} Fix a basis $\theta_1,\dots,\theta_n$ of $K/\Q$, let $p_\theta\colon T_{G^n}\to T_G$,
$(x_j)\mapsto\sum\theta_jx_j$, and $\Phi=\Psi\circ p_\theta$. Write the coordinates of $T_{G^n}$ as
$z_i,w_i,z_\rho,w_\rho,a_1,a_2\in\C^n$ and let $U\subset G^n$ be the vector subgroup
$\{z_i=z_\rho=0,\ \theta^{\mathsf T}a_1=0,\ \theta^{\mathsf T}(w_i+w_\rho+a_2)=0\}$, of dimension $4n-2$, defined over $\Qb$, with
$T_U\subset\ker\Phi$. Then $G'':=G^n/U=\Ga\times G_\ell$, where $G_\ell$ is the extension of $A=E_i^n\times E_\rho^n$ by
$\Ga$ obtained from the universal vector extension of $A$ by pushing out along $(w_i,w_\rho)\mapsto\theta^{\mathsf T}w_i+\theta^{\mathsf T}w_\rho$.
In the coordinates $s_1=\theta^{\mathsf T}a_1$, $z_i$, $z_\rho$, $s_2=\theta^{\mathsf T}(w_i+w_\rho+a_2)$ the map induced by
$\Phi$ is
\[
\Psi''(s_1,z_i,z_\rho,s_2)=\bigl(\theta^{\mathsf T}z_i+\theta^{\mathsf T}z_\rho+s_1,\;s_2\bigr).
\]
We apply \cite[Thm.~4.1]{W88} (Theorem~\ref{thm:W88}) to $G''$ with $d=2n+2$, $d_0=1$ ($\pi_0=s_1$), $d_1=0$,
$G_2=G_\ell$, $d_2=2n+1$; $V''=\Psi''^{-1}(V)$, a hyperplane, so $N=2n+1$; $W''=\ker\Psi''$, rational of dimension
$2n$; and $Y''$ generated by the images of $3n$ elements of $\Lambda(G^n)$ whose $\Phi$-images are $D\theta_hw_1$,
$D\theta_hw_2$ (taken in $\Omega(G^n)$, $D$ a common denominator) and $\theta_hv$, $h=1,\dots,n$. Since $U$ is a vector
group, $\Omega(G'')$ is the image of $\Omega(G^n)$, so $\Psi''(\Omega(G''))\subset P$. By Step~1 with $L=0$,
$\kappa=\mathrm{rk}(Y''\cap\Omega(G''))=2n$, and $\Gamma=\exp Y''$ has rank $n$. The inequality of
Theorem~\ref{thm:W88} becomes
\begin{equation}\label{eq:star}
\lambda+2\delta_2\le(\delta-\tau)(2n+2),
\end{equation}
and, with $r=\dim\Psi''(T_{G'})$, one has $\tau=\delta+r-2$; so an obstruction must have $r\le1$.

\emph{Step 3: the subgroups of $G''$.} Abelian subvarieties of $A$ are $B=B_i\times B_\rho$ ($E_i,E_\rho$ are
non-isogenous), and $B_i\subset E_i^n$ corresponds to a $\Q(i)$-subspace $S_i\subset\Q(i)^n$ with $T_{B_i}=\C\,\sigma(S_i)$,
$\dim B_i=\dim S_i$; likewise for $E_\rho$. With $\psi,\phi$ as in Lemma~\ref{lem:K} (for $k=\Q(i)$ and for $k=\Q(\rho)$), $\ker\psi\cap\ker\phi=0$ and each
kernel has $k$-dimension $n/2$. Because $[a]$ acts on $T(E^\natural)$ by $\mathrm{diag}(a,\bar a)$ and $\mathrm{Hom}(E^\natural,\Ga)=0$ (the universal vector extension is anti-affine; \cite{MM74}),
the image in $G_\ell$ of the universal extension of $B$ is a splitting $s(B)\cong B$ (with tangent space
$T_B\times0$) if $\phi(S_i)=\phi(S_\rho)=0$, and is $\pi^{-1}(B)$, a non-split extension, otherwise (a splitting would give a nonzero homomorphism from a quotient of $(E_i^\natural)^{b_i}\times(E_\rho^\natural)^{b_\rho}$ to $\Ga$). A connected subgroup over $B$ either contains the fibre or is a splitting, and a splitting is unique since $\mathrm{Hom}(B,\Ga)=0$. It follows that the
connected algebraic subgroups of $G''=\Ga\times G_\ell$ are exactly
\begin{center}
\begin{tabular}{llll}
type & $G'$ & condition & $\Psi''(T_{G'})$\\\hline
A & $\Ga\times\pi^{-1}(B)$ & --- & $\C^2$\\
B & $\Ga\times s(B)$ & $\phi(S)=0$ & $\C(1,0)$\\
C & $0\times\pi^{-1}(B)$ & --- & $\C(0,1)$ if $\psi(S)=0$, else $\C^2$\\
D & $0\times s(B)$ & $\phi(S)=0$ & $0$ if $B=0$, else $\C(1,0)$\\
E & $\{(\mu\,\mathrm{pr}_F(g),g):g\in s(B)\times F\}$, $\mu\ne0$ & $\phi(S)=0$ & $\C(\mu,1)$ if $B=0$, else $\C^2$
\end{tabular}
\end{center}
(with $\phi(S)=0$ meaning $\phi(S_i)=\phi(S_\rho)=0$, and similarly for $\psi$). In type E, $F\cong\Ga$ is the fibre of
$G_\ell\to A$ (coordinate $s_2$), $\pi^{-1}(B)=s(B)\times F$, and $\mathrm{pr}_F$ is the projection to $F$; so $G'$ is the
graph of $\mu\,\mathrm{pr}_F$ into the first factor $\Ga$ (coordinate $s_1$), and for $B=0$ it is $\{s_1=\mu s_2\}$. In every case allowed by the conditions,
$b=\dim B\le n$.

\emph{Step 4: no obstruction.} If $\Psi''(T_{G'})\in\{0,\C(1,0),\C(0,1)\}$ then $\lambda=n$: if
$\exp(\sum c_h\tilde y_h)\in G'$ for the lifts $\tilde y_h$ of $\theta_hv$, then $xv\in L+P$ with $x=\sum c_h\theta_h\in K$, so
$x=0$ by Step~1. With $\delta_0=0$ for types A, B, E and $\delta_0=1$ for C, D, inequality (\ref{eq:star}) reads:
$5n+2\le4n+4$ for $G'=0$; $5n+2-2b\le2n+2$ for types B and D ($B\neq0$); $5n-2b\le2n+2$ for type C with $r=1$; and
$\lambda+4n+2\le2n+2$ for type E with $B=0$. Since $b\le n$, none of these holds for $n\ge3$, and types with $r=2$ fail
$\delta>\tau$. So no connected $G'\ne G''$ satisfies the conclusion of Theorem~\ref{thm:W88}, a contradiction. \qed

\begin{remark}[The count]\label{rem:count}
Theorem~\ref{thm:RW} predicts Proposition~\ref{prop:D}: with $\kappa=2$ and $\alpha_K(G)=2(2+2/n)$
(Lemma~\ref{lem:natural}) it gives $\dim_K(V\cap\Psi(\Lambda_K(G)))\le2+4/n$, which is $<3$ only for $n>4$; with the
weight $\alpha(G'')/n=(4n+2)/n$ one gets $2+2/n<3$ for $n\ge3$, which is the threshold of Step~4.

A second derivation of Proposition~\ref{prop:D} applies Theorem~\ref{thm:W88k} to $(G'',V'',W'',Y'')$: $d=2n+2$,
$d_0=1$, $d_1=0$, $d_2=2n+1$, $t=\dim W''=2n$, $m=3n$, $\kappa=2n$, and (5.4) reads $3n>(4n+2-2n)\cdot1$, i.e.\
$3n>2n+2$ (for $K=\Q(i,\sqrt3)$, $n=4$: $12>18-\kappa=10$). It remains to see that no connected $G'\ne0$ has
$T_{G'}\subset V''$. First, $T_{G'}\not\subset\ker\Psi''$. If $G'$ meets the kernel $\Ga\times F$ of $G''\to A$ in a
subgroup of positive dimension, $T_{G'}$ contains a vector $(c,0,0,d)\ne0$, whose image is $(c,d)$. Otherwise $G'$ is
finite over its image $B$ in $A$, and $B\ne0$. Then $G'$ is an abelian variety, so $G'\subset0\times G_\ell$
($\mathrm{Hom}(B,\Ga)=0$) is a splitting $s(B)$; by Step~3 this forces $\phi(S)=0$, and then $\psi\ne0$ on $S$
(Lemma~\ref{lem:K}), so $\Psi''(T_{G'})=\C(1,0)$. Since $T_{G'}$ is $\Qb$-rational, it contains a $\Qb$-rational $x$ with
$\Psi''(x)\ne0$. If $T_{G'}\subset V''$, then $\Psi''(x)\in V$ is a nonzero $\Qb$-rational point of $V=\C w_1$, so
$\vp/(\beta\eta_\rho)\in\Qb$. This contradicts $\vp/\eta_\rho\notin\Qb$, which is W\"ustholz's theorem \cite{Wu89} (a
period of $E_i$ and a quasi-period of $E_\rho$ are $\Qb$-linearly independent). The same exclusion can be read off the
table in Step~3. Unlike Step~4, this derivation uses $\vp/\eta_\rho\notin\Qb$, and it rests on Theorem~\ref{thm:W88k} as
published.
\end{remark}

\smallskip
Corollary~\ref{cor:MX1} is the case $X=(u,\zeta_i(u))$, $Y=0$ (the logarithm of the origin), resp.\ $X=0$; the value $0$
of the ratio is excluded separately, since $1,u,\zeta(u)$ are $\Qb$-linearly independent for non-torsion $u$ on a CM
curve \cite{La80}; it also follows from \cite{Wu89} applied on $E^\natural\times\Ga$. Note that the base and fibre coordinates of a single logarithm cannot be chosen independently:
$w(X)-\zeta^*(z(X))\in\Qb$.

\begin{remark}[Other pairs of curves]\label{rem:general}
Let $E,E'$ be non-isogenous CM curves with lattices $\omega\Z[\alpha]$, $\omega'\Z[\alpha']$ (up to finite index), so that
$(\alpha-\bar\alpha)\,\omega\,\eta^*(\omega)=2\pi i=(\alpha'-\bar\alpha')\,\omega'\eta^*(\omega')$. With
$V=\C(\omega,\beta\eta^*(\omega'))$ the vectors $w_1=(\omega,\beta\eta^*(\omega'))$ and
$w_2=\bigl(\omega',\beta\frac{\alpha-\bar\alpha}{\alpha'-\bar\alpha'}\eta^*(\omega)\bigr)$ are parallel, and the same proof gives: if
$X\in\Lambda(E^\natural)$, $Y\in\Lambda(E'^\natural)$ and $(z(X)+z(Y)+a)\beta\eta^*(\omega')=(w(X)+w(Y)+t)\omega$, then $a=0$
and $z(X)\in k\omega$, $z(Y)\in k'\omega'$.
\end{remark}

\begin{remark}\label{rem:gamma}
By the explicit formulas for $\vp$ and $\Omega_\rho$, $\vp\Omega_\rho/\pi\in\Qb^\times\cdot\Gamma(1/4)^2\Gamma(1/3)^3/\pi^{5/2}$, and
$\vp/\eta_\rho=\tfrac{\sqrt3}{2\pi}\vp\Omega_\rho$. If $\vp/\eta_\rho=c\in\Qb$, then $X=Y=0$, $a=1$, $t=\beta/c$ satisfy the
hypothesis of Theorem~\ref{thm:MX} with $a\ne0$; so Corollary~\ref{cor:gamma} is the case $X=Y=0$. It also follows
directly from W\"ustholz's theorem \cite{Wu89} (a period of $E_i$ and a quasi-period of $E_\rho$ are $\Qb$-linearly
independent).
\end{remark}

\subsection{Controls}
(1) If either curve has no CM there is no relation making two periods parallel, so $\kappa$ drops to $n$ in Step~2 and
$G'=0$ satisfies the inequality of Theorem~\ref{thm:W88} ($5n+2\le6n+4$): no conclusion. The same happens if only one
Legendre relation is used. If the fibre action were $\mathrm{diag}(a,a)$ instead of $\mathrm{diag}(a,\bar a)$, subgroups
of type D with image $0$ would satisfy it; and for a rational line $V=\C(\mu,1)$ with true data on it ($\kappa=0$) the
type E line does. So the conjugate character, the second Legendre relation and the irrationality of $V$ are all used. (2) With the same curve twice, the antidiagonal of $E^\natural\times E^\natural$ lies in $\ker\Psi$
and the count collapses to the one-curve cell, which only gives $(\zeta(u)+t)/(u+a)\neq\pi/\vp^2$ --- a weak form of
Chudnovsky's theorem \cite{Ch84}. (3) Data on the period plane satisfy the relation only trivially.

\subsection{A twisted variant}\label{sec:MXP}
Exchanging the roles of base and fibre on $E_\rho$ gives a second statement of the same kind. Let
$G=E_i^\natural\times E_\rho^\natural\times\Ga^2$ and
\[
\Psi'(X_i,X_\rho,a_1,a_2)=\bigl(z(X_i)+w(X_\rho)+a_1,\;w(X_i)+z(X_\rho)+a_2\bigr),\qquad V'=\C\,w_1',\quad
w_1'=(\vp,\beta\Omega_\rho).
\]
With $w_2'=(\eta_\rho,\tfrac{2\beta}{\sqrt3}\eta_i)$ one has $w_2'=(\eta_\rho/\vp)\,w_1'$ by the two Legendre relations, and for
$K\supseteq\Q(i,\sqrt{-3},\beta)$ both lie in $\Psi'(\Omega_K(G))=(K\vp+K\eta_\rho)\times(K\eta_i+K\Omega_\rho)$.

\begin{theorem}[Theorem MX$'$]\label{thm:MXP}
Let $X\in\Lambda(E_i^\natural)$, $Y\in\Lambda(E_\rho^\natural)$, $a,b\in\Qb$ and $\beta\in\Qb^\times$. If
\[
\bigl(z(X)+w(Y)+a\bigr)\cdot\beta\,\Omega_\rho=\bigl(w(X)+z(Y)+b\bigr)\cdot\vp ,
\]
then $z(X)=z(Y)=0$ and both sides vanish.
\end{theorem}

\begin{proof}[Proof (sketch)]
\emph{Proposition D$'$: for $K\supseteq\Q(i,\sqrt{-3},\beta)$ of degree $n\ge8$, $V'\cap\Psi'(\Lambda_K(G))=Kw_1'+Kw_2'$.}
With $\theta$, $p_\theta$ and $\Phi'=\Psi'\circ p_\theta$ as before, a canonical lift $B^\natural$ of an abelian
subvariety $B_{S_i}\times B_{S_\rho}$ has tangent space in $\ker\Phi'$ only if $\psi(S_\bullet)=\phi(S_\bullet)=0$, i.e.\
$B=0$ (Lemma~\ref{lem:K}). So the largest connected subgroup $U'$ with $T_{U'}\subset\ker\Phi'$ is the vector subgroup
$\{\theta^{\mathsf T}(w_\rho+a_1)=0,\ \theta^{\mathsf T}(w_i+a_2)=0\}$, and $G''=G^n/U'$ is the pushout of
$(E_i^n\times E_\rho^n)^\natural$ along $(w_i,w_\rho)\mapsto(\theta^{\mathsf T}w_i,\theta^{\mathsf T}w_\rho)$, an extension of
$E_i^n\times E_\rho^n$ by $\Ga^2$ of dimension $2n+2$. It is a quotient of an anti-affine group, so
$\mathrm{Hom}(G'',\Ga)=\mathrm{Hom}(G'',\Gm)=0$: $d_0=d_1=0$, $d_2=2n+2$. With the lifts of $D\theta_hw_1'$,
$D\theta_hw_2'$ and $\theta_hv$ for a hypothetical $v\notin Kw_1'+Kw_2'$ one has $m=3n$ and $\kappa=2n$. The
$\kappa$-refined Theorem~1.1 of \cite[\S5d]{W88} (Theorem~\ref{thm:W88k}) applies: $m>(d_1+2d_2-\kappa)(d-1-t)$ reads $3n>2n+4$, i.e.\ $n>4$,
hence $n\ge8$ since $4\mid n$. Its remaining hypothesis holds because $\Psi''(T_{G'})$ would be a $\Qb$-rational subspace of
$V'$, which has no $\Qb$-rational point ($\Omega_\rho/\vp\notin\Qb$), so $T_{G'}\subset\ker\Psi''$ and $G'=0$ by the
maximality of $U'$. (A full case check against Theorem~\ref{thm:W88}, with the analogue of Step~1 of the proof of
Theorem~\ref{thm:MX}, gives the same threshold.)

\emph{Conclusion.} By D$'$, $z(X)+w(Y)+a=x\vp+x'\eta_\rho$ and $w(X)+z(Y)+b=x\beta\Omega_\rho+x'\tfrac{2\beta}{\sqrt3}\eta_i$
with $x,x'\in K$. The analytic subgroup theorem gives $z(X)=p\vp$, $z(Y)=q\Omega_\rho$ with $p\in\Q(i)$, $q\in\Q(\rho)$;
then $\zeta_i(p\vp)\in\bar p\,\eta_i+\Qb$ and $\zeta_\rho(q\Omega_\rho)\in\bar q\,\eta_\rho+\Qb$, and the $\Qb$-independence of
$1,\vp,\eta_\rho$ and of $1,\eta_i,\Omega_\rho$ gives $x=p$, $x'=\bar q$, $q=p\beta$, $\bar p=\tfrac{2\beta}{\sqrt3}\bar q$, and
vanishing constant terms. If $p\ne0$ this gives $\sqrt3|p|^2=2|q|^2$ with $|p|^2,|q|^2\in\Q_{>0}$, which is impossible.
\end{proof}

\begin{corollary}\label{cor:MXP}
For every nonzero $u\in\Lg_{E_i}$ (off the lattice) and $a,b\in\Qb$:
$(\zeta_i(u)+b)/(u+a)\notin\Qb^\times\cdot\Omega_\rho/\vp=\Qb^\times\cdot\Gamma(1/3)^3/(\sqrt\pi\,\Gamma(1/4)^2)$.
\end{corollary}

\noindent Together with Corollary~\ref{cor:MX1}, the ratio $(\zeta_i(u)+t)/(u+a)$ avoids both $\Qb\,\pi/(\vp\Omega_\rho)$ and
$\Qb^\times\Omega_\rho/\vp$.

\subsection{Two degree-two period statements}\label{sec:deg2}
Propositions~\ref{prop:D} and~D$'$ also apply to data that are periods plus algebraic constants. This gives statements
about products of periods of the two curves.

\begin{corollary}\label{cor:deg2}
\begin{itemize}
\item[(i)] The numbers $\pi,\ \vp,\ \Omega_\rho,\ \vp\Omega_\rho$ are linearly independent over $\Qb$; equivalently, so are
$\pi,\eta_i,\eta_\rho,\eta_i\eta_\rho$.
\item[(ii)] The numbers $1,\ \vp,\ \Omega_\rho,\ \vp\Omega_\rho/\pi$ are linearly independent over $\Qb$; equivalently, so are
$\pi,\vp,\eta_\rho,\vp\eta_\rho$.
\end{itemize}
\end{corollary}

\begin{proof}
The equivalences follow by multiplying a relation by $\eta_i\eta_\rho/\pi$, resp.\ by $\eta_\rho$, and using
$\vp\eta_i=\pi$, $\Omega_\rho\eta_\rho=2\pi/\sqrt3$. We use that $1,\pi,\vp,\eta_i,\Omega_\rho,\eta_\rho$ are linearly
independent over $\Qb$ (analytic subgroup theorem \cite{Wu89} on $E_i^\natural\times E_\rho^\natural\times\Gm\times\Ga$).

(i) A relation with zero coefficient on $\vp\Omega_\rho$ is excluded by that independence, so suppose
$\pi=A\vp\Omega_\rho+B\vp+C\Omega_\rho$ with $A\ne0$. Take $\beta=A$, $X=(\vp,\eta_i)\in\Lambda(E_i^\natural)$, $Y=0$, $a=C/A$,
$b=-B$ in Theorem~\ref{thm:MXP}: the left side is $(\vp+C/A)A\Omega_\rho=\pi-B\vp$ and the right side is
$(\eta_i-B)\vp=\pi-B\vp$. So the hypothesis holds, and Theorem~\ref{thm:MXP} gives $z(X)=0$, whereas $z(X)=\vp$.

(ii) Suppose $\vp\Omega_\rho/\pi=a+b\vp+c\Omega_\rho$ (a relation with zero left-hand coefficient is excluded by the
independence above). If $c=0$, dividing by $\vp$ and using $\pi/\vp=\eta_i$, $\Omega_\rho=2\pi/(\sqrt3\eta_\rho)$ gives
$\Omega_\rho=a\eta_i+b\pi$. If $b=0$, multiplying by $\eta_\rho$ gives $\tfrac{2}{\sqrt3}\vp=a\eta_\rho+\tfrac{2c}{\sqrt3}\pi$. Both
contradict the independence above, so $b,c\ne0$. Put $t=-1/c$, $\beta=-\sqrt3\,b/(2c)$, $a_1=a/b$, and take
$X=(\vp,\eta_i+t)\in\Lambda(E_i^\natural)$, $Y=0$. Then $v=(\vp+a_1,\eta_i+t)\in\Psi(\Lambda(G))$ lies on
$V=\C(\vp,\beta\eta_\rho)$: expanding $(\vp+a_1)\beta\eta_\rho=(\eta_i+t)\vp$ with $\vp\eta_i=\pi$ and
$\eta_\rho=2\pi/(\sqrt3\Omega_\rho)$ gives exactly the assumed relation. Proposition~\ref{prop:D} gives
$v=x(\vp,\beta\eta_\rho)+x'(\Omega_\rho,\tfrac{2\beta}{\sqrt3}\eta_i)$ with $x,x'\in K$. The first coordinate and the
independence of $1,\vp,\Omega_\rho$ give $a_1=0$, $x=1$, $x'=0$; then $\eta_i+t=\beta\eta_\rho$, contradicting the
independence of $1,\eta_i,\eta_\rho$.
\end{proof}

\noindent Both statements are genuinely of degree two: rewriting with the Legendre relations always leaves a product
of periods of the two different curves, so they are not consequences of the analytic subgroup theorem. They are
predicted by the period conjecture (by Chowla--Selberg \cite{Gr78} they concern $\pi$, $\Gamma(1/4)$, $\Gamma(1/3)$, whose algebraic
independence is open), and they do not follow from known algebraic independence results. They have the shape of a
nonvanishing $2\times2$ determinant of periods, like the theorem of Barr\'e-Sirieix, Diaz, Gramain and Philibert \cite{BDGP96} in the form
\cite[Cor.~41]{W06}. We
found them nowhere in print; they are the most delicate claims of the paper and we flag them for an expert check.
Neither statement contains the constant $1$ together with $\vp\Omega_\rho$, so they do not give $\vp\Omega_\rho\notin\Qb$. For
$E_i,E_\rho$, that statement is the CM case of Waldschmidt's conjecture \cite[slide~2]{W03s} that $\omega_1\omega_2^*-\omega_1^*\omega_2$ is
transcendental for non-isogenous curves, since $\vp\cdot\rho\Omega_\rho-\Omega_\rho\cdot i\vp=(\rho-i)\vp\Omega_\rho$. We do not
claim it.

\subsection{The abelian-surface analogue}\label{sec:MXA}
A simple CM abelian surface $A$ has two eigen-pairs of periods with $P_1P_4,\ P_2P_3\in\pi\,\Qb^\times$ (Riemann's bilinear
relations). For $A=\mathrm{Jac}(y^2=x^5+1)$, $P_a=(1-\zeta_5^a)B(\tfrac a5,\tfrac12-\tfrac a5)/5$, and the reflection
formula gives these relations exactly. With $G=A^\natural\times\Ga^2$, $\Psi=(z_1+z_2+a_1,\,w_3+w_4+a_2)$ and
$V=\C(P_1,\beta P_3)$, the same argument (with $\alpha_K(A^\natural)=4+4/n$) shows: for non-torsion $P\in A(\Qb)$ with
logarithm $v$, $(z_1(v)+z_2(v)+a)\beta P_3\neq(w_3(v)+w_4(v)+t)P_1$. We state this only as a proof sketch; the steps not
written out are: the computation $\alpha_K(A^\natural)=4+4/n$ (no abelian part in $H$, since the four coordinates carry
four distinct characters of $\Q(\zeta_5)$); condition (i) via W\"ustholz on $A^\natural$, which uses that $A^\natural$ has
no proper algebraic subgroup surjecting onto $A$; the description of $\Lambda(A^\natural)$; and the final deduction that
$v$ is torsion. This sketch uses Theorem~\ref{thm:RW}. The direct argument of \S\ref{sec:MX} should carry over with the same
numbers ($G''=\Ga\times G_\ell$ of dimension $2n+2$, abelian subvarieties given by $\Q(\zeta_5)$-subspaces and four
conditions from the four embeddings), but we have not written it out.

\section{Single-partner limits for $e^{\pi^2}$}\label{sec:3E}

The three exponentials conjecture asks that $\lambda_1\lambda_2=\gamma\lambda_3$ with $\lambda_j\in\Lg$, $\gamma\in\Qb^\times$
force $\lambda_3=0$; $e^{\pi^2}\notin\Qb$ is the instance $\lambda_1=\lambda_2=i\pi$. In the count above, a configuration
of rank-one shape gives a theorem only if $\kappa$ exceeds what the true data need. For lines with a single
coordinate per factor, $\kappa\le1$, and one CM factor supplies exactly one unit of weight. On the period side,
Laurent's identity (\ref{eq:laurent}) and its abelian analogue show that these configurations are linearisable:
multiplying by $\pi$ turns them into linear relations among $1$-periods, decided by the analytic subgroup theorem.

To get $\kappa=2$ on a line one needs two independent periods with parallel images, i.e.\ a true \emph{homogeneous}
degree-two relation $p_ap_d=p_bp_c$ among periods. The CM Legendre relations $(\alpha-\bar\alpha)\,\omega\,\eta^*(\omega)=2\pi i$
of two non-isogenous CM curves provide one; so do two eigen-pairs of a CM abelian variety. The multiplicative group
cannot take part: its only period is $2\pi i$, and the corresponding relation would pair $2\pi i$ with $1$, which is not a
period (for $E^\natural\times\Gm$ with $\Gm$ added to either coordinate, parallel periods force a $\Qb$-linear relation
among $1,\omega,\eta^*$, so $\kappa\le1$). Under the hypothesis $e^{\pi^2}\in\Qb$ the number $\pi^2$ is a logarithm,
not a period, so it adds data, not $\kappa$. This is why Theorem~\ref{thm:MX} exists while the same count cannot reach
$e^{\pi^2}$. The $\Gm$-analogue of Theorem~\ref{thm:MX} (replace $\Omega_\rho$ by $i\pi$ and the $E_i$-data by logarithms)
is $i\pi\cdot\log\beta=\gamma\log\alpha$, an open case of the conjecture.

With two partners the mechanism does reach $e^{\pi^2}$: Corollary~\ref{cor:MG1} (one of $e^{\pi^2}$, $e^{\pi\Omega}$,
$e^{\pi^2/\Omega}$) uses $\pi^2$ and $\pi^2/\Omega$ as data and the period $(2\pi i,2ic\Omega)$ for $\kappa$. In the
configurations we enumerated (lines in $\C^2$, at most three factors per coordinate; not refereed), no single-partner
disjunction ``one of $e^{\pi^2}$, $X$'' appears beyond (M3) and the known strong-six partners $e^{c/\pi}$, $e^{c\pi^3}$ (from
the strong six exponentials theorem \cite[Cor.~2]{Roy92} with $x=(1,\pi)$ and $y=(1/\pi,1,\pi)$, resp.\ $y=(1,\pi,\pi^2)$).

Within the monomial model (heuristic; not refereed), a slack count explains this (proof sketch). On a line $V=\C(1,r)$,
each slot contributes monomials: its period, plus the constant $1$ for $\Ga$ and for the fibre. So $\Ga$ gives $\{1\}$,
$\Gm$ gives $\{\pi\}$, the base of $E$ gives $\{\Omega\}$, and the fibre of $E^\natural$ gives $\{\pi/\Omega,1\}$. Their weights in the
count are $0$, $1$, $1$, $1+2/n$, and the true relations are homogeneous for the grading $\deg\pi=2$,
$\deg(\text{1-period})=1$. A line count is then a matching between the monomials of the two coordinates. The hypothesis
$e^{\pi^2}\in\Qb$ adds the monomial $\pi^2$ (degree $4$) to the $\Gm$-slots, and a partner hypothesis adds one monomial
$m_X$. A single-partner
disjunction needs the partner to sit opposite $\pi^2$, so $m_X=\pi^2r^{\pm1}$, or the $e^{\pi^2}$-edges alone to fill the
slack. Running through the cases leaves exactly $m_X=\pi\cdot p$ with $p$ a period of the slot carrying the partner, i.e.\
$\wp(c\pi\Omega)$ (the case (M3)), together with the strong-six partners $e^{c/\pi}$ and $e^{c\pi^3}$. An exhaustive
enumeration in a bounded range (at most four slots per coordinate) confirms this. On lines of type MX the slope has
degree $0$, so a $\Gm$-slot can never be matched; $e^{\pi^2}$ can enter such a count only together with a second hypothesis.

\section{Consistency checks}\label{sec:checks}
The main results are consequences of a published theorem (completed by Lemma~\ref{lem:U} in the case used) applied to
explicit groups, so an error would most likely sit in the bookkeeping. We checked that the same bookkeeping does not
prove false statements, and that it gives no conclusion in the control configurations below, where it should not.
\begin{itemize}
\item \emph{Rational lines.} If $V$ contains a $\Qb$-rational point, the statement is false (e.g.\ $\Gm$ in place of $E$,
$s=2$, $x=\log2$). In each configuration we checked (MG, MG$^\natural$, MX) the count finds an obstruction, typically a
subtorus or a $\Ga$-line; for MG and MG$^\natural$ it occurs at equality in the inequality of Theorem~\ref{thm:W88}.
\item \emph{No CM, wrong field, one Legendre relation, wrong coordinates.} An obstruction appears (the trivial subgroup,
or a type-D subgroup) in each of these cases: a non-CM curve; $k\not\subset K$; only one Legendre relation on an MX-type
line; for MX, $\mathrm{diag}(a,a)$ in place of $\mathrm{diag}(a,\bar a)$ on $T(E^\natural)$; for MG$^\natural$, a non-eigen fibre
coordinate. So no conclusion is claimed in these cases. (For MG$^\natural$ the conjugation itself is not load-bearing;
what matters is that $\Psi$ reads an eigen-coordinate.)
\item \emph{Pure periods.} Data consisting of periods alone are true, so they test the period term by itself. In all the
configurations we built (semi-abelian $G$ with $V$ a hyperplane containing a corank-one sublattice of periods,
$\Gm^{d_1}\times E^g$, the groups of \S\S\ref{sec:MG}, \ref{sec:MX}), an obstruction exists.
\item \emph{Sharpness of the period term.} For the graph of complex multiplication in $E\times E$ ($\kappa=2$), the
diagonal in $\Gm^2$ ($\kappa=1$), and the Weil-restricted graph of complex multiplication with maximal $\kappa=n$, the
obstruction occurs exactly at equality. A period term larger by one would prove false statements there. The
applications in \S\S\ref{sec:MG}, \ref{sec:MX} use $\kappa=\mathrm{rk}(V''\cap\ker\exp)$ exactly. This is calibration evidence,
not a proof of \cite[Prop.~6.1]{W88}.
\item \emph{$\Gm$ alone} [proved]. For $G=\Gm^r$ the period space $\Psi(\Omega_K)$ is $2\pi i$ times a $K$-rational
subspace. So a line $V=\C(1,s)$ meets it only if $s\in K$, where the statement is false and the method is silent. With a
logarithm in place of the elliptic period ($\kappa=0$), the same count reproduces the six exponentials theorem (a
calibration).
\item \emph{Special exponents.} Values of $t$ that make one power algebraic ($2^t=3$, $2^t=1$) are not known to lie
outside $\Qb^\times\Omega/\pi$. They lie inside only under degree-two relations such as $\pi\log3\in\Qb\,\Omega\log2$ or
$\Omega\log2\in\Qb\pi^2$, which are conjecturally false. If they did lie inside, the resulting consequences of
Corollary~\ref{cor:MG2} would be conjecturally true and would contradict no known result.
\end{itemize}

\section{Status and sources not consulted}\label{sec:status}

Theorems~\ref{thm:MG}, \ref{thm:MGn}, \ref{thm:MX} and~\ref{thm:MXP} (hence Theorem~\ref{thm:A} and the corollaries)
rest on the following.
\begin{itemize}
\item Theorem~\ref{thm:W88} \cite[Thm.~4.1]{W88}. Its proof \cite[\S\S6--7]{W88} uses:
\begin{itemize}
\item the auxiliary function with periods and the $\kappa$-saving \cite[Prop.~2.1, 2.4]{W83};
\item multiplicities along $W$ (Lemma~\ref{lem:U}; a proof sketch, checked by two independent internal referees, not
peer-reviewed). Its published inputs are \cite[Prop.~2.1, Lemma~2.2]{W83}, Serre's embedding and height bound
\cite{Se79}, Bertrand's \cite[App.~I, Lemme~7 and its Corollaire, with Lemme~1]{W79} (cited in \cite[p.~396]{W88} as
``Lemma~7''), and the lower bound of \cite[p.~640]{W83};
\item Philippon's zero estimate \cite{Ph86};
\item a complete system of addition laws of one fixed bidegree for the projective embedding of $G_2$. This exists for
any quasi-projective commutative group variety by an elementary argument (Remark~\ref{rem:addition}); for abelian
varieties see also Lange--Ruppert \cite{LR85}. (Serre's appendix \cite{Se79} gives the compactification with its
$G$-action, but not addition laws; cf.\ also \cite{FW84}, which we have not checked.)
\end{itemize}
\item For the sketch of Theorem~\ref{thm:MXP} and the second derivations (Remarks~\ref{rem:MG5d}, \ref{rem:MGn5d},
\ref{rem:count}): Theorem~\ref{thm:W88k}, as published in \cite[\S5d]{W88}.
\item Legendre's relation; W\"ustholz's analytic subgroup theorem \cite{Wu89}; Masser \cite{Ma75}; Baker;
the transcendence of $\Omega$ and of $\Omega/\pi$ (Schneider, Masser
\cite{Ma75}).
\item Standard facts on CM abelian varieties and universal vector extensions.
\end{itemize}
For $u\in\Lg_E$ (all of Corollary~\ref{cor:M}), Theorem~\ref{thm:M} follows from Laurent's identity and the analytic
subgroup theorem; for $u\in\Qb\Lg_E$ it rests on Statement R (in the hyperplane case, which is Theorem~\ref{thm:W88k}
in substance).
Extensions to CM abelian varieties (four exponentials at Beta values such as $\Gamma(1/5)^2/\Gamma(2/5)$) and a mechanism using
first-kind period relations of CM abelian varieties, as well as strong four exponentials at all CM slopes
$\Qb^\times\{\Omega^{\pm1},(\Omega/\pi)^{\pm1}\}$, are treated separately. What most needs an expert read: Lemma~\ref{lem:U} (the one step of \cite[Prop.~6.1]{W88} that is only indicated in
print); the subgroup classifications in \S\S\ref{sec:MG} and~\ref{sec:MX}; and the degree-two statements of
Corollary~\ref{cor:deg2}, which say more about $\pi$, $\Gamma(1/4)$, $\Gamma(1/3)$ than any known algebraic-independence
result. Corollary~\ref{cor:deg2}(i) rests on Theorem~\ref{thm:MXP}, whose write-up has had one internal check.
The claim that MX is not a
linearisation, and the discussion in \S\ref{sec:3E}, are heuristic (they use the period conjecture). Sources we
could not consult, any of which might contain these statements: Bertrand \cite{Be83} (the source of the $K$-structure used
in \cite{RW93}); Waldschmidt \cite{W90} (the general form of Theorem~\ref{thm:W88}); Baker--W\"ustholz \cite[\S6.2]{BW07};
Bertrand--Laurent \cite{BL81}; Faltings--W\"ustholz \cite{FW84};
Roy \cite{R1}; MathSciNet
citation searches. Numerical values above were computed with mpmath at 50--60 digits; they illustrate and prove nothing.

\appendix
\section{Multiplicities along $W$ in the auxiliary function}\label{app:U}

We use the notation of \cite[\S6]{W88}. $G=\Ga^{d_0}\times\Gm^{d_1}\times G_2$ over a number field $K_0$ with $d_2\ge1$ and
$G_2$ embedded in $\mathbb P_M$ by Serre's embedding (coordinates $\psi_0,\dots,\psi_M$, entire of order $\le2$;
$\xi=\psi/\psi_0$). $V\subset T_G$ has dimension $N<d$, $W\subset V$ is $K_0$-rational of dimension $t$, $y_1,\dots,y_m\in V$
with $\exp_G(y_j)\in G(K_0)$, $Y(S)=\{\sum h_jy_j: 0\le h_j\le S\}$, $0\le\kappa\le\mathrm{rk}_\Z(V\cap\ker\exp_G)$, and
$\Delta^{d-N}=C\,S^{d_1+2d_2-\kappa}(\log S)^{bd_0}$, $T(\log S)^a=D_0(\log S)^b=D_1S=D_2S^2=\Delta$.

\begin{lemma}\label{lem:U}
Assume $2N\ge d_1+2d_0+\kappa$, $d_2\ge1$, $a,b\ge1$. For $\nu$ large and $S\ge S_0(\nu)$ there is a polynomial
$P_S\in K_0[U,V,W]$ of degree $\le D_0$ in $U$, $\le D_1$ in $V$, homogeneous of degree $\le D_2$ in $W$, not vanishing
identically on $G$, which vanishes with multiplicity $\ge T$ along $W$ at every point of $\exp_G(Y(S))$. Hence
\cite[Prop.~6.1]{W88} holds under these assumptions.
\end{lemma}

\begin{proof}[Sketch]
\emph{Step 1} is \cite[Prop.~2.1 and p.~640]{W83} unchanged. With $r=\nu^2S$, $R=\nu^3S$, it gives a nonzero
$Q_S\in\Z[U,V,W]$ of degrees $\le D_i/\nu$ and coefficients $\le e^{\Delta/\nu}$ with $|Q_S(\lambda,\mu,\psi)(z)|\le
e^{-\nu^6\Delta}$ for $z\in V$, $|z|\le\nu^2S$. The hypothesis $2N\ge d_1+2d_0+\kappa$ is the condition
$U\ge r^2\log(R/r)$ of Prop.~2.1 (printed with $-\kappa$ on p.~641, a misprint). No condition along $W$ enters the linear
system, so the saving $r^{\kappa}$ in the hypothesis of \cite[Prop.~2.1, p.~633]{W83} is untouched. We also use the lower bound
$\max_i|\psi_i(t+z)|\ge e^{-c(\nu S)^2}$ for $t\in Y(\nu S)\cap\mathcal T$ and $|z|\le\rho$, where $Y=\mathcal T+F$ with $F$
finite (\cite[p.~640]{W83}, \cite[Lemme~1.2.2]{W79}; the $\psi_i$ have no common zero \cite[Prop.~8]{Se79}). In Step~2 it
bounds the scalar $\mu_j(z)$ in $E^{(j)}(\xi(\gamma),\xi(z))=\mu_j(z)\psi(t+z)$. (The polynomial variables $U,V,W$ below
are not the subspaces $V,W$.)

\emph{Step 2} (derivatives). Fix $t\in Y(\nu S)\cap\mathcal T$, $\gamma=\exp_G(t)$, and an addition law $E^{(j)}$ of the
embedding with $E^{(j)}(\xi(\gamma),\xi(0))\ne0$. For $z=\sum s_iw_i\in W$ small put
$F_t(s)=Q_S(\lambda(\gamma)+\lambda(z),\mu(\gamma)\mu(z),E^{(j)}(\xi(\gamma),\xi(z)))$. By homogeneity, $F_t$ is a unit
times $Q_S(\lambda,\mu,\psi)(t+z)/\psi_0(t+z)^{D^\circ}$, so it has the same order along $W$. Its Taylor coefficients
$\theta'$ of order $<T$ satisfy $|\theta'|\le e^{-\nu^6\Delta+c\nu\Delta}$ (Cauchy). They are polynomials, with no
denominator depending on $\gamma$, in the coefficients of $Q_S$, the coordinates of $\gamma$, and the Taylor coefficients
of $\lambda,\mu,\xi$ along $W$ at the origin. The latter are $K_0$-rational with house and denominator
$\le c^{|\tau|+1}(|\tau|!)^{c}$: by Bertrand's Corollaire to his Lemme~7 \cite[App.~I, p.~177]{W79} the coordinates of the
exponential map have finite arithmetic functional order, every point being regular, and his Lemme~1 \cite[App.~I,
pp.~166--167]{W79} turns this into such bounds for derivatives (\cite[p.~396]{W88} cites this as ``Lemma~7''). Since
$|\tau|<T$ and $T\log T\le c\Delta$ (as $a\ge1$ and $\log\Delta\le c\log S$), Serre's height bounds for $\gamma$ then give
$h(\theta')\le c\nu\Delta$ with $c$ independent of $\nu$. Liouville's inequality forces $\theta'=0$ once $\nu^5$ exceeds a constant. (Dividing by a
power of $E_0^{(j)}(\xi(\gamma),\xi(0))$ instead would cost a height $T(\nu S)^2$, which is too large.)

\emph{Step 3} (translation). Write $Y(S)\subset(Y(\nu S)\cap\mathcal T)+F'$ with $F'$ finite. For $f\in F'$ take the
product over the addition laws $E^{(j)}$ with $E^{(j)}(\cdot,\xi(-\exp f))\not\equiv0$ of $Q_S$ translated by $-\exp f$;
then take the product over $f\in F'$. Translations commute with invariant derivations, and each factor is a unit times
the translate on a dense open set, so the product vanishes with multiplicity $\ge T$ along $W$ on $\exp_G(Y(S))$. It is
not identically zero on $G$, and its degrees are $\le D_i$ once $\nu$ exceeds (number of laws)$\cdot$(bidegree)$\cdot|F'|$.
No norm to $\Q$ is taken, since \cite[Prop.~6.1]{W88} allows $\Qb$-coefficients.

The exponent $b\ge1$ changes only the number of unknowns, which is compensated exactly by the factor $(\log S)^{bd_0}$ of
$\Delta^{d-N}$, and lowers the height term $(D_0/\nu)\log S$. In the proof of \cite[Thm.~4.1]{W88}, $b>1$ is used only in
the case $\delta=\tau$, which requires $d_0\ge1$ (Theorem~\ref{thm:MX}).
\end{proof}

\begin{remark}[Addition laws]\label{rem:addition}
Let $G_2\subset\mathbb P_M$ be locally closed over $K_0$, and $X=G_2\times G_2\subset\mathbb P_M\times\mathbb P_M$. The local
ring of $X$ at a point $p$ is a quotient of that of $\mathbb P_M\times\mathbb P_M$, whose elements are ratios of
bihomogeneous forms of equal bidegree. Clearing denominators, near $p$ the group law $X\to G_2\subset\mathbb P_M$ is given
by bihomogeneous forms $(A_0,\dots,A_M)$ of one bidegree, not all zero at $p$. By quasi-compactness finitely many such
systems suffice, and multiplying each by suitable monomials $x_i^ay_l^b$ gives a complete system of addition laws of one
bidegree $(\beta,\beta)$, with coefficients in a finite extension of $K_0$. (This replaces $K_0$ by a finite extension in
Step~2, which only changes $[K_0:\Q]$ in the Liouville inequality.) Step~2 uses these laws only on
$G_2(K_0)\times\exp(\text{a neighbourhood of }0\text{ in }W)$, and Step~3 only with a fixed second argument.
\end{remark}

\begin{thebibliography}{99}
\bibitem{Abe20} Y. Abe, A transcendence result on elliptic functions using functions of two variables,
arXiv:2005.01211 (2020; withdrawn 2021).
\bibitem{BDGP96} K. Barr\'e-Sirieix, G. Diaz, F. Gramain, G. Philibert, Une preuve de la conjecture de Mahler--Manin,
Invent.\ Math.\ 124 (1996), 1--9.
\bibitem{BL81} D. Bertrand, M. Laurent, Propri\'et\'es de transcendance de nombres li\'es aux fonctions th\^eta, C. R.
Acad.\ Sci.\ Paris S\'er.\ I 292 (1981), 747--749.
\bibitem{ChCh26} G. V. Chudnovsky, D. V. Chudnovsky, Irrationality exponents of real CM elliptic periods and powers of the
lemniscate period, arXiv:2610.31112 (2026), with Lean files.
\bibitem{Ch84} G. V. Chudnovsky, \emph{Contributions to the Theory of Transcendental Numbers}, Math.\ Surveys Monogr.\ 19,
Amer.\ Math.\ Soc., 1984.
\bibitem{Di04} G. Diaz, Utilisation de la conjugaison complexe dans l'\'etude de la transcendance de valeurs de la
fonction exponentielle usuelle, J. Th\'eor.\ Nombres Bordeaux 16 (2004), 535--553.
\bibitem{Gr78} B. H. Gross, On the periods of abelian integrals and a formula of Chowla and Selberg, with an appendix by
D. E. Rohrlich, Invent.\ Math.\ 45 (1978), 193--211.
\bibitem{FW84} G. Faltings, G. W\"ustholz, Einbettungen kommutativer algebraischer Gruppen und einige ihrer
Eigenschaften, J. reine angew.\ Math.\ 354 (1984), 175--205.
\bibitem{LR85} H. Lange, W. Ruppert, Complete systems of addition laws on abelian varieties, Invent.\ Math.\ 79
(1985), 603--610.
\bibitem{Ma76} D. W. Masser, A note on a paper of Franklin, Acta Arith.\ 31 (1976), 143--152.
\bibitem{Ph86} P. Philippon, Lemmes de z\'eros dans les groupes alg\'ebriques commutatifs, Bull.\ Soc.\ Math.\ France 114
(1986), 355--383.
\bibitem{BK77} W. D. Brownawell, K. K. Kubota, The algebraic independence of Weierstrass functions and some related
numbers, Acta Arith.\ 33 (1977), 111--149.
\bibitem{BW07} A. Baker, G. W\"ustholz, \emph{Logarithmic Forms and Diophantine Geometry}, New Math.\ Monogr.\ 9,
Cambridge Univ.\ Press, 2007.
\bibitem{Be83} D. Bertrand, Endomorphismes de groupes alg\'ebriques; applications arithm\'etiques, in \emph{Approximations
diophantiennes et nombres transcendants} (Luminy 1982), Progr.\ Math.\ 31, Birkh\"auser, 1983, 1--45.
\bibitem{La80} M. Laurent, Ind\'ependance lin\'eaire de valeurs de fonctions doublement quasi-p\'eriodiques, C. R. Acad.\
Sci.\ Paris S\'er.~A 290 (1980), 397--399.
\bibitem{La82} M. Laurent, Transcendance de p\'eriodes d'int\'egrales elliptiques II, J. reine angew.\ Math.\ 333 (1982),
144--161.
\bibitem{Ma75} D. W. Masser, \emph{Elliptic Functions and Transcendence}, Lecture Notes in Math.\ 437, Springer, 1975.
\bibitem{R1} D. Roy, Transcendance et questions de r\'epartition dans les groupes alg\'ebriques, in \emph{Approximations
diophantiennes et nombres transcendants} (Luminy 1990), de Gruyter, 1992, 249--274.
\bibitem{Roy92} D. Roy, Matrices whose coefficients are linear forms in logarithms, J. Number Theory 41 (1992), 22--47.
\bibitem{RW93} D. Roy, M. Waldschmidt, Autour du th\'eor\`eme du sous-groupe alg\'ebrique, Canad.\ Math.\ Bull.\ 36
(1993), 358--367.
\bibitem{W88} M. Waldschmidt, On the transcendence methods of Gel'fond and Schneider in several variables, in \emph{New
Advances in Transcendence Theory} (Durham 1986), Cambridge Univ.\ Press, 1988, 375--398.
\bibitem{MM74} B. Mazur, W. Messing, \emph{Universal Extensions and One Dimensional Crystalline Cohomology}, Lecture
Notes in Math.\ 370, Springer, 1974.
\bibitem{SMF84} D. Bertrand et al., Les nombres transcendants, M\'em.\ Soc.\ Math.\ France (2) 13 (1984).
\bibitem{Se79} J.-P. Serre, Quelques propri\'et\'es des groupes alg\'ebriques commutatifs, Appendice~II in \cite{W79},
191--202.
\bibitem{W79} M. Waldschmidt, \emph{Nombres transcendants et groupes alg\'ebriques}, Ast\'erisque 69--70 (1979), with
appendices by D. Bertrand and J.-P. Serre.
\bibitem{W83} M. Waldschmidt, Sous-groupes analytiques de groupes alg\'ebriques, Ann.\ of Math.\ 117 (1983), 627--657.
\bibitem{W05b} M. Waldschmidt, Further variations on the six exponentials theorem, Hardy-Ramanujan J.\ 28 (2005), 1--9.
\bibitem{W05} M. Waldschmidt, Variations on the six exponentials theorem, in \emph{Algebra and Number Theory}
(Hyderabad 2003), Hindustan Book Agency, 2005, 338--355.
\bibitem{W06} M. Waldschmidt, Elliptic functions and transcendence, in \emph{Surveys in Number Theory}, Dev.\ Math.\ 17,
Springer, 2008, 143--188.
\bibitem{W07} M. Waldschmidt, The role of complex conjugation in transcendental number theory, in \emph{Diophantine
Equations} (DION 2005, Mumbai), ed.\ N. Saradha, Tata Inst.\ Fund.\ Res.\ Stud.\ Math.\ 20, Narosa, New Delhi, 2008,
297--308.
\bibitem{W22} M. Waldschmidt, The four exponentials problem and Schanuel's conjecture, in \emph{Mathematics Going
Forward}, Lecture Notes in Math.\ 2313, Springer, Cham, 2023, 579--592.
\bibitem{W90} M. Waldschmidt, Dependence of logarithms of algebraic points, in \emph{Number Theory, Vol.~II} (Budapest
1987), Colloq.\ Math.\ Soc.\ J\'anos Bolyai 51, North-Holland, 1990, 1013--1035.
\bibitem{W03s} M. Waldschmidt, Dependence of logarithms on commutative algebraic groups, slides, Hyderabad, December 2003.
\bibitem{W19} M. Waldschmidt, Third kind elliptic integrals and transcendence, appendix to C. Bertolin, Third kind
elliptic integrals and 1-motives, J. Pure Appl.\ Algebra 224 (2020), no.~10, 106396; arXiv:1905.07247.
\bibitem{WW85} J. Wolfart, G. W\"ustholz, Der \"Uberlagerungsradius gewisser algebraischer Kurven und die Werte der
Betafunktion an rationalen Stellen, Math.\ Ann.\ 273 (1985), 1--15.
\bibitem{Wu89} G. W\"ustholz, Algebraische Punkte auf analytischen Untergruppen algebraischer Gruppen, Ann.\ of Math.\ 129
(1989), 501--517.
\end{thebibliography}

\end{document}
